\chapter{JavaScript e TypeScript}
\thispagestyle{empty} 
\label{cap:js-ts} 


Este capítulo apresenta uma introdução prática e objetiva às linguagens \textit{JavaScript} (JS) e \textit{TypeScript} (TS), destacando suas principais características e como elas são aplicadas no desenvolvimento Web. O conteúdo é dividido em três partes para facilitar a compreensão:

\begin{itemize}
  \item \textbf{Visão Geral:} Aqui, você conhecerá os fundamentos de \textit{JavaScript} e \textit{TypeScript}, suas diferenças e vantagens. Esses conceitos são abordados nas Seções~\ref{JS_visaoGeral} e~\ref{JS_visaoGeralTS}.

  \item \textbf{Sintaxe e Funcionalidades:} Esta parte cobre a estrutura das linguagens de forma prática. Serão abordados tipos de dados, variáveis, operações matemáticas e manipulação de \textit{strings} (Seções~\ref{JS_funcionamento} a \ref{JS_operacoesBasica}). Em seguida, exploramos estruturas essenciais como condicionais (\texttt{if}/\texttt{else}, \texttt{switch}/\texttt{case}) e laços de repetição (\texttt{for}, \texttt{while}) de forma integrada (Seções~\ref{JS_condicionais} e~\ref{JS_lacosRepeticao}).   Também veremos funções e suas variações, como \textit{arrow functions} (Seção~\ref{JS_funcoes}), além do tratamento de exceções e a programação assíncrona com \textit{callbacks}, \textit{Promises} e \texttt{async/await} (Seções~\ref{JS_tratamentoExcecoes},~\ref{JS_promises} e~\ref{JS_callback}). Para organização de dados, abordamos \textit{arrays} e coleções chaveadas (Seções~\ref{JS_arays} e ~\ref{JS_colecoesChaveadas}). Por fim, introduzimos classes, herança e interfaces (Seções~\ref{JS_classesObjetos} e~\ref{JS_interfaces}).

  \item \textbf{Aplicação Prática:} Para consolidar os conceitos, a Seção~\ref{JS_applicacaoExemplo} apresenta um exemplo completo: a implementação de um jogo \textit{Flappy Bird} utilizando \textit{JavaScript} e \textit{TypeScript}.
\end{itemize}


%Finalmente, a Seção~\ref{JS_tratamentoExcecoes} discute o tratamento de exceções, seguido pela exploração das \textit{Promises} na Seção~\ref{JS_promises} e o uso de \textit{callbacks} na Seção~\ref{JS_callback}. As seções~\ref{JS_arays} e ~\ref{JS_colecoesChaveadas} abordam \textit{arrays} e coleções chaveadas, respectivamente, enquanto a Seção~\ref{JS_classesObjetos} explora o conceito de classes e objetos, incluindo a herança de classes. A Seção~\ref{JS_interfaces} aborda o conceito de interfac


%Este capítulo apresenta uma introdução às linguagens \textit{JavaScript} (JS) e \textit{TypeScript} (TS), destacando suas características e aplicações no desenvolvimento Web. Inicialmente, a Seção~\ref{JS_visaoGeral} apresenta uma visão geral sobre JavaScript, abordando sua tipagem dinâmica, programação assíncrona e ecossistema. Em seguida, a Seção~\ref{JS_visaoGeralTS} fornece uma introdução ao TypeScript, destacando os benefícios da tipagem estática e explorando o ecossistema relacionado.


%Posteriormente,  Seção~\ref{JS_funcionamento} explica o funcionamento das linguagens, enquanto a Seção~\ref{JS_tipoDadosVariaveis} detalha os tipos de dados e variáveis, incluindo a declaração, constantes, variáveis e estruturas de dados. As Seções~\ref{JS_numeros} e~\ref{JS_string} discutem números e \textit{strings}, respectivamente. Já a Seção~\ref{JS_operacoesBasica} explora as operações básicas, como atribuição, comentários, operações aritméticas, uso do objeto \texttt{Math}, operações relacionais e entrada/saída. Em seguida, a Seção~\ref{JS_condicionais} aborda o uso de condicionais, com ênfase nas estruturas \texttt{if}/\texttt{else} e \texttt{switch}/\texttt{case}, enquanto a Seção~\ref{JS_lacosRepeticao} explora os laços de repetição for e \texttt{while}. A Seção~\ref{JS_funcoes}, apresenta conceitos fundamentais sobre funções, incluindo a introdução da \textit{arrow function} como uma sintaxe mais concisa. 

%Finalmente, a Seção~\ref{JS_tratamentoExcecoes} discute o tratamento de exceções, seguido pela exploração das \textit{Promises} na Seção~\ref{JS_promises} e o uso de \textit{callbacks} na Seção~\ref{JS_callback}. As seções~\ref{JS_arays} e ~\ref{JS_colecoesChaveadas} abordam \textit{arrays} e coleções chaveadas, respectivamente, enquanto a Seção~\ref{JS_classesObjetos} explora o conceito de classes e objetos, incluindo a herança de classes. A Seção~\ref{JS_interfaces} aborda o conceito de interface e a Seção~\ref{JS_applicacaoExemplo} demostra um exemplo de uso do \textit{JavaScript} e do \textit{TypeScript}. Ao abordar todos esses tópicos de forma paralela para ambas as linguagens, essa seção visa estabelecer uma base sólida de conhecimento teórico necessário para compreender os aspectos práticos e as aplicações posteriores discutidas no trabalho.


\section{Visão Geral sobre \textit{JavaScript}}
\label{JS_visaoGeral}

O \textit{JavaScript} é uma linguagem de programação de alto nível e multiparadigma, suportando programação orientada a objetos, funcional e imperativa incluído orientada a objetos, amplamente utilizada no desenvolvimento de aplicações Web. Originalmente criada para fornecer interatividade aos sites, permitindo a manipulação dinâmica do conteúdo da página e a resposta a eventos do usuário, o \textit{JavaScript} evoluiu ao longo dos anos e se tornou uma linguagem versátil, com suporte a programação assíncrona e desenvolvimento tanto no cliente quanto no servidor~\cite{mdn-basics}.

A história do \textit{JavaScript} remonta à década de 1990, quando foi criado por Brendan Eich na Netscape Communications Corporation. Inicialmente chamada de \textit{LiveScript}, a linguagem foi renomeada para \textit{JavaScript} para capitalizar a popularidade da linguagem \textit{Java} na época. Em 1997, a Ecma International lançou a especificação ECMAScript, que é a padronização oficial da linguagem. Desde então, várias versões do ECMAScript foram lançadas, introduzindo melhorias e novos recursos para o \textit{JavaScript}~\cite{w3schools-history, javascript_mdn}.

O \textit{JavaScript} é uma linguagem interpretada, o que significa que o código fonte é executado diretamente por um interpretador em tempo de execução, sem a necessidade de compilação prévia. Os navegadores modernos incorporam motores \textit{JavaScript}, como o V8, \textit{SpiderMonkey} e \textit{JavaScriptCore}, para executar o código \textit{JavaScript}~\cite{javascript_mdn}. Na \autoref{fig:ex-exec-js} é possível visualizar um exemplo de um código JavaScript implementado a uma página HTML sendo interpretado por um navegador Web.

\begin{figure}
    \centering
    \includegraphics{figs/cap_js/ex-exec-js.pdf}
    \caption{Código JavaScript sendo interpretado pelo navegador Web}
    \label{fig:ex-exec-js}
\end{figure}

Uma característica importante do \textit{JavaScript} é a tipagem dinâmica, o que dispensa a necessidade de declaração explícita de tipos de variáveis. Variáveis em \textit{JavaScript} podem armazenar valores de diferentes tipos, como números, strings, objetos e funções. Além disso, o \textit{JavaScript} suporta programação assíncrona por meio de recursos como chamadas de retorno (\textit{callbacks}), \textit{Promises} e a adição mais recente do async/await, que permitem a execução de tarefas demoradas em segundo plano, sem bloquear a execução do programa~\cite{devpleno, mdn-dynamic-typing, mdn-data-structures}.

O \textit{JavaScript} desempenha um papel fundamental no desenvolvimento Web e se tornou uma linguagem padrão para páginas da Web. Sua importância reside no fato de que o \textit{JavaScript} é responsável por tornar as páginas Web dinâmicas, interativas e funcionais. Assim, é possível adicionar recursos como validação de formulários, animações, manipulação do \textit{DOM} (Modelo de Objeto de Documento) e interações em tempo real com o usuário~\cite{javascript-rockcontent}.

Além do desenvolvimento \textit{front-end} de sites e aplicações Web, o \textit{JavaScript} é amplamente adotado em \textit{frameworks} e bibliotecas populares, como \textit{React.js}, \textit{Angular.js},  \textit{Vue.js} e \textit{Node.js}. Esses frameworks possibilitam a construção de aplicativos Web complexos, tanto no lado do cliente quanto no servidor, além de fornecerem uma estrutura robusta para o desenvolvimento de interfaces de usuário sofisticadas e responsivas. O \textit{JavaScript} também é amplamente utilizado no desenvolvimento de jogos, aplicativos móveis híbridos e até no desenvolvimento de servidores com o \textit{Node.js}~\cite{javascript_mdn}.

O \textit{Node.js} é uma versão do \textit{JavaScript} que se destaca por sua aplicação no desenvolvimento \textit{back-end}. Enquanto o \textit{JavaScript} é amplamente reconhecido como a linguagem de programação padrão para páginas Web, responsável por torná-las dinâmicas e interativas, o \textit{Node.js} estende o uso dessa linguagem para o lado do servidor. Essa expansão permitiu que os desenvolvedores utilizassem o \textit{JavaScript} em todas as camadas de uma aplicação, unificando a linguagem e facilitando o compartilhamento de código entre o \textit{front-end} e o \textit{back-end}.

\subsection{Tipagem Dinâmica}

Uma das características mais distintivas do \textit{JavaScript} é a tipagem dinâmica. Diferentemente de linguagens com tipagem estática, como \textit{C++} ou \textit{Java}, em que os tipos de variáveis devem ser declarados explicitamente, o \textit{JavaScript} permite que as variáveis armazenem valores de diferentes tipos em momentos distintos durante a execução do programa. Essa flexibilidade da tipagem dinâmica pode ser uma vantagem, pois simplifica o desenvolvimento, permitindo que variáveis sejam reutilizadas para diferentes propósitos sem a necessidade de redeclaração ou conversões explícitas de tipo. Por exemplo, uma variável pode armazenar um número em um momento e, em seguida, armazenar uma string ou objeto posteriormente.

Em contrapartida, a tipagem dinâmica também pode apresentar desafios, especialmente em projetos grandes e complexos, nos quais o controle estrito dos tipos pode ser desejado para evitar erros. É importante ter cuidado ao lidar com tipos em \textit{JavaScript} para garantir que as operações sejam realizadas corretamente e que os valores sejam coerentes com o esperado. E como solução a esta problemática com a linguagem, surgiu o \textit{TypeScript}.

\subsection{Programação Assíncrona}

Outra característica importante do \textit{JavaScript} é o suporte à programação assíncrona. Isso permite que o código seja executado de forma não sequencial, em que tarefas demoradas ou bloqueantes podem ser executadas em segundo plano, enquanto outras operações continuam. Antes dos recursos assíncronos do \textit{JavaScript}, as chamadas de função que envolviam operações demoradas, como solicitações de rede, poderiam bloquear a execução do programa até que a operação fosse concluída. Isso poderia levar a uma experiência ruim para o usuário, pois a interface do usuário poderia parecer congelada até que a operação fosse concluída.

No entanto, com a introdução de recursos como chamadas de retorno (\textit{callbacks}), \textit{Promises} e async/await, o \textit{JavaScript} oferece mecanismos poderosos para lidar com operações assíncronas. As chamadas de retorno permitem que uma função seja passada como argumento para outra função, a ser executada quando uma determinada operação for concluída. As \textit{Promises} são objetos que representam o resultado futuro de uma operação assíncrona, permitindo que o código se organize de forma mais legível e evitando a ocorrência de ``\textit{callback hell}'' (aninhamento excessivo de chamadas de retorno). O \textit{async/await} é uma adição mais recente à linguagem que simplifica ainda mais a programação assíncrona, permitindo que o código assíncrono seja escrito de forma mais similar à programação síncrona, facilitando a compreensão e a manutenção do código.

Esses recursos de programação assíncrona são especialmente úteis no contexto de aplicações Web, já que muitas vezes é necessário realizar operações demoradas, como fazer solicitações a APIs externas, buscar dados em bancos de dados ou processar grandes volumes de informações. O \textit{JavaScript} oferece ferramentas poderosas para lidar com essas situações, permitindo que as aplicações sejam responsivas e mantenham uma boa experiência para o usuário.

\subsection{Ecossistema do \textit{JavaScript}}

O \textit{JavaScript} possui um ecossistema vibrante, com uma grande quantidade de bibliotecas, \textit{frameworks} e ferramentas disponíveis para facilitar o desenvolvimento de aplicações Web. Além disso, o \textit{JavaScript} é suportado em todos os principais navegadores modernos, tornando-se uma linguagem universalmente aceita para o desenvolvimento \textit{front-end}. Existem muitos \textit{frameworks} populares construídos com \textit{JavaScript}, que fornecem estruturas e padrões para o desenvolvimento Web. Alguns dos mais conhecidos~são:

\begin{itemize}
\item \textbf{\textit{React.js}\footnote{\textbf{React.js:} https://react.dev/}}: Um \textit{framework} \textit{JavaScript} mantido pelo Facebook, que é amplamente utilizado para construir interfaces de usuário interativas. O \textit{React.js} utiliza um modelo de componentes reutilizáveis, que simplifica o desenvolvimento e a manutenção de interfaces complexas.

\item \textbf{\textit{Angular.js}\footnote{\textbf{Angular.js:} https://angular.io/}}: Um \textit{framework} \textit{JavaScript} mantido pelo Google, que permite a criação de aplicações Web escaláveis e de alta performance. O \textit{Angular.js} utiliza o conceito de \textit{data binding}, em que as alterações nos dados são automaticamente refletidas na interface do usuário, simplificando a sincronização entre a lógica e a apresentação.

\item \textbf{ \textit{Vue.js}\footnote{\textbf{Vue.js:} https://vuejs.org/}}: Um \textit{framework} \textit{JavaScript} progressivo e de fácil aprendizado, que permite a construção de interfaces de usuário interativas e reativas. O \textit{Vue.js} oferece uma sintaxe intuitiva e flexível, permitindo que os desenvolvedores criem componentes reutilizáveis e construam aplicações de forma incremental.

\item \textbf{ \textit{Node.js}\footnote{\textbf{Node.js:} https://nodejs.org/en}}: Uma plataforma de tempo de execução \textit{JavaScript} construída sobre o mecanismo V8 do Google Chrome. O \textit{Node.js} permite a execução de código \textit{JavaScript} no lado do servidor, o que abre um mundo de possibilidades para o desenvolvimento Web. Com o \textit{Node.js}, é possível criar servidores Web, APIs RESTful, microsserviços e muito mais.
\end{itemize}

Esses são apenas alguns exemplos do vasto ecossistema \textit{JavaScript}. Existem inúmeras outras bibliotecas e ferramentas disponíveis, cada uma com suas características e casos de uso específicos. O \textit{JavaScript} continua evoluindo rapidamente, com novas versões do ECMAScript sendo lançadas regularmente, trazendo melhorias e recursos adicionais para a linguagem.


\section{Visão Geral sobre \textit{TypeScript} \label{JS_visaoGeralTS}}

O \textit{TypeScript} é uma linguagem de programação desenvolvida pela Microsoft que se baseia no \textit{JavaScript} padrão, adicionando recursos de tipagem estática e outros recursos avançados. Foi projetada para solucionar algumas das limitações do \textit{JavaScript}, oferecendo aos desenvolvedores uma experiência de desenvolvimento mais robusta e segura. Assim como o \textit{JavaScript}, o \textit{TypeScript} são essenciais para desenvolvimento Web, cada um com vantagens únicas para diferentes cenários, tanto no \textit{front-end} quanto no \textit{back-end}. Ele compartilha a mesma sintaxe básica e estrutura de programação do \textit{JavaScript}, mas introduz a adição de tipos estáticos opcionais. Essa adição permite que os desenvolvedores especifiquem o tipo de cada variável, parâmetro e retorno de função, trazendo benefícios em termos de detecção de erros em tempo de compilação e ferramentas de autocompletar mais poderosas.

O \textit{TypeScript} foi lançado como um projeto de código aberto em 2012, com base na especificação \textit{ECMAScript}, a mesma base do \textit{JavaScript}. Ele transpila para \textit{JavaScript}, o que significa que o código escrito em \textit{TypeScript} é traduzido para \textit{JavaScript} para ser executado em qualquer ambiente compatível com \textit{JavaScript}. O \textit{TypeScript} é amplamente adotado em projetos de grande escala, nos quais a manutenção e a escalabilidade são essenciais. \textit{Frameworks} populares, como \textit{Angular.js}, adotam o \textit{TypeScript} como sua linguagem principal, aproveitando os recursos adicionais fornecidos pela linguagem para facilitar o desenvolvimento de aplicativos Web robustos.

A \autoref{fig:comp-ts-to-js} ilustra um exemplo de transpilação de um código \textit{TypeScript} em um código \textit{JavaScript}. No código do arquivo \texttt{exemplo.ts} é possível observar a declaração de uma função \texttt{somar}, que especifica os tipos dos parâmetros da função, sendo eles dois valores do tipo \texttt{number}, e a função, por sua vez, retorna a soma desses valores. Já no código do arquivo \texttt{exemplo.js}, resultado da traspilaço do \textit{TypeScript}, apresenta o mesmo código, contudo, sem os tipos estáticos definidos. Além disso, o compilador acrescenta na primeira linha de código o comando ``\texttt{use strict}'', que impede que variáveis sejam usadas sem antes serem declaradas em \textit{JavaScript}~\cite{w3schools_js_strict}.

\begin{figure}
    \centering
    \includegraphics[width=0.5\textwidth]{figs/cap_js/comp-ts-to-js.pdf}
    \caption{Exemplo de Compilação de JavaScript para TypeScript}
    \label{fig:comp-ts-to-js}
\end{figure}

Uma das principais vantagens do \textit{TypeScript} é a detecção de erros em tempo de transpilação. Ao adicionar tipos estáticos ao código, o compilador do \textit{TypeScript} pode identificar erros, como atribuição incorreta de tipos, chamadas de função inválidas e propriedades inexistentes. Essa verificação estática ajuda a reduzir erros comuns e permite que os desenvolvedores identifiquem problemas antes mesmo de executar o código. Além da tipagem estática, o \textit{TypeScript} oferece recursos avançados, como interfaces, herança de classe, classes abstratas, módulos e uma variedade de recursos adicionais para facilitar o desenvolvimento de aplicativos complexos. Ele também suporta as últimas atualizações do ECMAScript, permitindo que os desenvolvedores utilizem as funcionalidades mais recentes do \textit{JavaScript}.



\subsection{Benefícios da Tipagem Estática}

A adição de tipos estáticos ao \textit{JavaScript} no \textit{TypeScript} traz diversos benefícios para o desenvolvimento de software. Alguns dos principais benefícios incluem:

\begin{itemize}
\item \textbf{Detecção de erros em tempo de transpilação}: Com a tipagem estática, o compilador do \textit{TypeScript} é capaz de identificar erros no código antes mesmo de ser executado. Isso ajuda a reduzir erros comuns, como atribuições incorretas de tipos e chamadas de função inválidas, proporcionando um código mais confiável e menos propenso a erros em tempo de execução.

\item \textbf{Melhor autocompletar e recursos de IDE}: A especificação dos tipos no \textit{TypeScript} permite que as IDEs forneçam recursos de autocompletar mais poderosos e informações sobre os tipos esperados. Isso aumenta a produtividade do desenvolvedor, reduzindo o tempo gasto em pesquisas de documentação e evitando erros de digitação.

\item \textbf{Refatoração segura}: Com a tipagem estática, é possível realizar refatorações de código com mais segurança. O compilador pode identificar e alertar sobre alterações que podem afetar outras partes do código, facilitando a manutenção e a evolução do projeto.

\item \textbf{Documentação mais clara}: A especificação dos tipos no código fonte torna a documentação do código mais clara e autoexplicativa. Os tipos fornecem informações valiosas sobre os parâmetros e retornos das funções, tornando o código mais legível e compreensível para outros desenvolvedores.

\item \textbf{Colaboração facilitada}: A tipagem estática ajuda na colaboração entre membros da equipe, especialmente em projetos grandes. Com os tipos definidos explicitamente, fica mais fácil entender as intenções do código e evitar conflitos entre diferentes partes do projeto.

\item \textbf{Melhor escalabilidade}: O \textit{TypeScript} facilita o desenvolvimento de projetos de grande escala de maneira mais eficiente. Com a tipagem estática, é possível ter uma visão mais clara das dependências e interações entre os diferentes componentes do sistema, facilitando a manutenção e a evolução do código à medida que o projeto cresce.
\end{itemize}

Desde 2012, o \textit{TypeScript} tem sido amplamente adotado por grandes corporações devido à sua tipagem estática e compatibilidade com JavaScript, manutenível e escalável. Esses benefícios são alcançados graças às características da linguagem, como a tipagem estática e os recursos de orientação a objetos, que permitem aos desenvolvedores detectar erros de forma antecipada, melhorar a clareza e legibilidade do código e facilitar sua manutenção ao longo do tempo. Além disso, o \textit{TypeScript} oferece suporte a ferramentas de desenvolvimento poderosas, como autocompletar e refatoração, que ajudam a acelerar o processo de desenvolvimento e garantir uma base de código mais robusta. Com essas vantagens, o \textit{TypeScript} se tornou a escolha preferida para muitas equipes de desenvolvimento que buscam construir sistemas mais confiáveis e escaláveis.

\subsection{Ecossistema \textit{TypeScript}}

Assim como o \textit{JavaScript}, o \textit{TypeScript} também possui um ecossistema ativo e diversificado, com uma ampla gama de bibliotecas e \textit{frameworks} disponíveis para facilitar o desenvolvimento de aplicações Web. Muitas bibliotecas e \textit{frameworks} populares do \textit{JavaScript} são compatíveis com o \textit{TypeScript}, permitindo que os desenvolvedores tirem proveito dos recursos adicionais de tipagem estática. Além disso, o \textit{TypeScript} é suportado por uma série de ferramentas de desenvolvimento, como IDEs (Ambientes de Desenvolvimento Integrado) e editores de código, que fornecem recursos avançados para ajudar os desenvolvedores a escrever, depurar e refatorar código \textit{TypeScript}. Algumas das ferramentas populares para desenvolvimento \textit{TypeScript} incluem:

\begin{itemize}
\item \textbf{\textit{Visual Studio Code}\footnote{ https://code.visualstudio.com/}}: Um editor de código gratuito e de código aberto desenvolvido pela Microsoft, com suporte nativo ao \textit{TypeScript}. O \textit{Visual Studio Code} oferece recursos avançados, como autocompletar, depuração integrada, integração com controle de versão e extensibilidade por meio de plugins.

\item \textbf{\textit{WebStorm}\footnote{ https://www.jetbrains.com/pt-br/webstorm/}}: Um IDE poderoso para o desenvolvimento Web, que oferece suporte ao \textit{TypeScript} e a outras tecnologias Web. O \textit{WebStorm} fornece recursos avançados de refatoração, depuração e análise de código, ajudando os desenvolvedores a escreverem código \textit{TypeScript} de forma eficiente.

\item \textbf{\textit{TypeScript Playground}\footnote{ https://www.typescriptlang.org/play}}: Uma ferramenta online fornecida pela equipe do \textit{TypeScript}, que permite experimentar e testar código \textit{TypeScript} diretamente no navegador. O \textit{TypeScript Playground} oferece recursos de autocompletar, verificação de erros em tempo real e visualização do \textit{JavaScript} gerado pelo código \textit{TypeScript}.

% \item \textbf{\textit{tslint} e \textit{eslint}}: Ferramentas de linting que ajudam a manter um código \textit{TypeScript} consistente e livre de erros comuns. Essas ferramentas analisam o código em busca de problemas, como estilo de código inconsistente, uso incorreto de recursos da linguagem e possíveis erros.

\item \textbf{\textit{Deno}\footnote{ https://deno.com/runtime}}: Um runtime \textit{JavaScript}/\textit{TypeScript} seguro e moderno, construído com base no mecanismo V8 do Google Chrome. O Deno permite executar código \textit{TypeScript} diretamente, sem a necessidade de compilação prévia, e oferece recursos avançados, como módulos nativos, controle de acesso a recursos do sistema e suporte a importação de módulos diretamente do navegador.
\end{itemize}

Esses são apenas alguns exemplos do vasto ecossistema \textit{TypeScript}. Existem inúmeras outras bibliotecas e ferramentas disponíveis, cada uma com suas características e casos de uso específicos. O \textit{TypeScript} continua evoluindo rapidamente, com novas versões sendo lançadas regularmente, trazendo melhorias e recursos adicionais para a linguagem.


\subsection{Instalando e Usando \textit{TypeScript}}

Para usufruir da tipagem estática e outros recursos avançados do \textit{TypeScript} no desenvolvimento do projeto é necessário instalá-lo, utilizando o \texttt{npm} (\textit{Node Package Manager}), que é o gerenciador de pacotes do Node.js. Para usar o \texttt{npm} basta baixar e instalar o Node.js, pois ele já vem incluído.

Para instalar o \textit{TypeScript} usando o \texttt{npm}, usa-se os comandos apresentados na \autoref{listing:instalar-ts}. O primeiro comando é utilizado para instalá-lo globalmente, já o segundo para a adicioná-lo ao projeto como dependência de desenvolvimento.

\begin{lstlisting}[language=bash, label=listing:instalar-ts, caption=Comandos de instalação do \textit{TypeScript}.]
npm install -g typescript
npm install --save-dev typescript
\end{lstlisting}

Para iniciar um projeto \textit{TypeScript} usa-se o comando apresentado da \autoref{listing:inicializacao-ts}. Esse comando gera um arquivo chamado \texttt{tsconfig.json}, onde pode-se definir configurações como o diretório de saída, nível de rigor e suporte a módulos.

\begin{lstlisting}[language=bash, label=listing:inicializacao-ts, caption=Comando de inicialização de projeto \textit{TypeScript}.]
npx tsc --init
\end{lstlisting}

Para transpilar arquivos \textit{JavaScript} (.js) para \textit{TypeScript} (.ts) utiliza-se o comando mostrado na \autoref{listing:transpilacao-ts}, onde \texttt{arquivo.ts} é o caminho do arquivo. Se o objetivo for transpilar todo o projeto basta omitir o caminho do arquivo.

\begin{lstlisting}[language=bash, label=listing:transpilacao-ts, caption=Comando de transpilação de arquivos \textit{TypeScript} (.ts) para \textit{JavaScript} (.js).]
npx tsc arquivo.ts
\end{lstlisting}

Utilizando Node.js, é possível executar arquivos \textit{TypeScript} sem precisar transpilá-los manualmente. Basta instalar o \texttt{ts-node} globalmente, ou localmente no projeto, como já apresentado na \autoref{listing:instalar-ts}, porém, ao invés de \texttt{typescript} é \texttt{node-ts}. Uma vez feito isso, pode-se usar o comando da \autoref{listing:exec-ts-node} para executar o arquivo em \textit{TypeScript}.

\begin{lstlisting}[language=bash, label=listing:exec-ts-node, caption=Execução direta de arquivo \textit{TypeScript}.]
ts-node arquivo.ts
\end{lstlisting}

\section{Funcionamento}
\label{JS_funcionamento}

O \textit{JavaScript} (JS) é uma linguagem interpretada, o que significa que os programas codificados nessa linguagem são executados por um interpretador em tempo de execução, sem a necessidade de compilação prévia para linguagem de máquina. Por outro lado, o \textit{TypeScript} é uma linguagem compilada, mas não para linguagem de máquina. Os arquivos \textit{TypeScript}, de extensão ``.ts'', passam por um processo de compilação para \textit{JavaScript}, gerando arquivos de extensão ``.js'', derivados dos arquivos \textit{TypeScript}. Somente então esses arquivos \textit{JavaScript} podem ser executados~\cite{cubos_academy, stxnext}.

O \textit{TypeScript} (TS) utiliza um compilador (tsc) para transformar o código \textit{TypeScript} em um código \textit{JavaScript} equivalente. Durante o processo de compilação, o compilador verifica a sintaxe do código TS, realiza análise estática de tipos e converte o código em \textit{JavaScript} válido. Esse código JS resultante pode ser executado em qualquer ambiente compatível com \textit{JavaScript}, como navegadores Web, servidores \textit{Node.js} e outros ambientes de tempo de execução~\cite{book_progtypescript}. 



\section{Tipos de Dados e Variáveis}
\label{JS_tipoDadosVariaveis}
Os tipos de dados referem-se a que tipo de informações uma linguagem consegue operar primitivamente, ou seja, o que é representável computacionalmente no contexto da linguagem de programação \cite{embarcados, wikipedia_tipodedado}. Já as variáveis, são espaços de memória alocados para armazenar esses tipos de dados. É por meio das variáveis que se manipulam os dados na programação~\cite{gaea, wikipedia_variavel}.

\subsection{Declaração}
Para o programa reconhecer uma variável, precisa-se declará-la, ou seja, dar o comando de alocar um espaço da memória do computador e nomeá-lo para que se possa armazenar e consultar dados nesse espaço de memória no decorrer do programa. O nome dado a esse trecho de memória é o nome da variável, também chamado de identificador~\cite{wikipedia_variavel, gaea}. 

Para declarar uma variável em \textit{JavaScript} e \textit{TypeScript} é bem simples, como mostra a \autoref{lst:dec-var}. Observa-se que o comando inicia com a palavra reservada \texttt{let}, que indica a declaração de uma variável, e logo em seguida, consta o nome da variável (não pode conter caracteres especiais). No caso do TS, ainda em sequência, têm-se a opção de inserir dois pontos e especificar o tipo de dado da variável declarada. De forma geral, pode-se substituir ``\texttt{minhaVariavel}'' por qualquer nome de variável (contanto que não contenha caracteres especiais ou comece com número) e ``\texttt{number}'' por qualquer tipo disponível na linguagem.

\begin{lstlisting}[language=TypeScript, label=lst:dec-var, caption=Exemplo de Declaração de Variável em JS e TS.]
let minhaVariavel;          // Em JavaScript   
let minhaVariavel: number;  // Em TypeScript
\end{lstlisting}



\subsection{Constantes e Variáveis}

A distinção entre constantes e variáveis é fundamental na programação. As constantes são valores imutáveis que não sofrem alterações ao longo do programa, fornecendo consistência e confiabilidade ao código. Essas constantes são frequentemente utilizadas para representar informações fixas, como números específicos, configurações ou valores que não devem ser modificados. Por outro lado, as variáveis são utilizadas para armazenar informações que podem ser alteradas durante a execução do programa. Elas permitem a flexibilidade de manipular e atualizar valores à medida que o programa avança. As variáveis são úteis quando se precisa lidar com dados que podem ser modificados, como contadores, resultados de cálculos ou entradas fornecidas pelo usuário. A distinção entre constantes e variáveis permite aos programadores controlar o comportamento do código de acordo com as necessidades do programa, garantindo assim um melhor gerenciamento dos dados e uma maior flexibilidade na manipulação dos valores.

A sintaxe geral para a declaração de constantes é ilustrada na \autoref{lst:dec-const}. No caso do \textit{JavaScript}, a declaração simplesmente consiste em utilizar a palavra \texttt{const} seguida pelo nome da constante. Já no \textit{TypeScript}, é possível também definir o tipo da constante através da notação: ``\texttt{<tipo>}'', onde \texttt{<tipo>} representa o tipo de dado que a constante irá armazenar. Essa adição de tipos é uma das vantagens do \textit{TypeScript}, pois traz maior segurança e ajuda a evitar erros de tipo durante o desenvolvimento.

\begin{lstlisting}[language=TypeScript, label=lst:dec-const , caption=Sintaxe Generalizada de Declaração de Constante em JS e TS.]
const <nomeDaConstante>          // em JavaScript
const <nomeDaConstante>: <tipo>  // em TypeScript
\end{lstlisting}

Uma outra maneira ainda de se declarar uma variável, como pode-se observar na \autoref{lst:dec-Var} é utilizando a palavra reservada \texttt{var}. A principal diferença entre o \texttt{let} e o \texttt{var} é que o \texttt{var} tem escopo de função ou escopo global, é içada (\textit{hoisted}) para o topo do escopo e permite reatribuição, enquanto o \texttt{let}, usado na \autoref{lst:dec-var}, tem escopo de bloco, não é içada (\textit{hoisted}) para o topo do escopo e também permite reatribuição, mas apenas dentro do mesmo bloco em que foi declarada \cite{statement-var, statement-let}. A preferência atual é usar \texttt{let} em vez de \texttt{var} na maioria dos casos, pois o escopo de bloco do \texttt{let} oferece uma melhor granularidade de controle e ajuda a evitar problemas relacionados ao escopo.

\begin{lstlisting}[language=TypeScript, caption=Sintaxe Generalizada de Declaração de Variável Global em JS., label=lst:dec-Var]
var <nomeDaVariavel>          // Em JavaScript
var <nomeDaVariavel>: <tipo>  // Em TypeScript
\end{lstlisting}

\subsection{Estruturas de Dados}

As variáveis em um algoritmo podem assumir diferentes tipos de dados, dependendo do papel que desempenham. O tipo de dados atribuído a uma variável define o conjunto de valores que ela pode armazenar e as operações que podem ser realizadas sobre ela~\cite{wikipedia_tipodedado}. Por exemplo, ao calcular a área de uma circunferência, é apropriado utilizar uma variável com tipo numérico para armazenar o raio, permitindo realizar operações matemáticas. Por outro lado, ao armazenar o nome de um usuário, é mais adequado utilizar uma variável com tipo textual para lidar com caracteres e strings. Essa seleção cuidadosa dos tipos de dados garante que as variáveis sejam usadas de forma correta e coerente dentro do contexto do algoritmo.

Além disso, nas linguagens de programação, existem os tipos de dados primitivos, que são os tipos mais básicos e fundamentais da linguagem. Esses tipos primitivos servem como base para a definição de outros tipos de dados mais complexos. Em geral, os tipos primitivos incluem inteiros, números de ponto flutuante, caracteres individuais, valores booleanos (verdadeiro/falso) e outros tipos simples. Eles possuem uma representação direta na memória e são processados de forma eficiente pelos sistemas computacionais. Os tipos primitivos fornecem os blocos de construção essenciais para a construção de estruturas de dados mais complexas e o desenvolvimento de algoritmos. Portanto, compreender e utilizar corretamente os tipos de dados primitivos é fundamental para a programação eficiente e eficaz~\cite{devmedia_tipoprimitivo, wikipedia_primitivedatatype}.

Os tipos de dados existentes em \textit{JavaScript} também são válidos em \textit{TypeScript}, mas vale lembrar que diferentemente da linguagem \textit{JavaScript}, \textit{TypeScript} é uma linguagem fortemente tipada e possui alguns tipos adicionais, como o \texttt{any} e o \texttt{never} \cite{vidafullstack, site_typescript, javascript_mdn}. Na \autoref{tabela_tipos} é possível relacionar os tipos e seus significados, como por exemplo: o tipo \texttt{Boolean} representa um valor lógico verdadeiro ou falso; o tipo \texttt{Null} representa a ausência intencional de qualquer objeto ou valor; e assim por diante.


\begin{table}
\centering
\caption{Tipos de dados em \textit{JavaScript} (JS) e \textit{TypeScript} (TS)}

\begin{tabular}{|p{2cm}|p{9cm}|}
\hline
\textbf{Tipo} & \textbf{Descrição} \\
\hline
\texttt{Boolean} & Representa um valor lógico verdadeiro ou falso. Pode ser \texttt{true} (verdadeiro) ou \texttt{false} (falso). \\
\hline
\texttt{Null} & Representa a ausência intencional de qualquer objeto ou valor. O \texttt{null} é um valor especial que indica a falta de um valor ou objeto. \\
\hline
\texttt{Undefined} & Representa uma variável que foi declarada, mas não recebeu um valor. Quando uma variável é declarada, mas não inicializada, ela terá o valor \texttt{undefined}. \\
\hline
\texttt{Number} & Representa valores numéricos, como inteiros e números de ponto flutuante. \\
\hline
\texttt{BigInt} & Representa valores inteiros maiores do que o limite superior do tipo \texttt{number}. \\
\hline
\texttt{String} & Representa uma sequência de caracteres. \\
\hline
\texttt{Symbol} & Representa um valor único e imutável que pode ser usado como chave de uma propriedade de objeto. \\
\hline
\texttt{Object} & Representa uma coleção de pares de chave/valor. \\
\hline
\end{tabular}
\label{tabela_tipos}
\end{table}

Em \textit{JavaScript} e \textit{TypeScript}, existem duas formas de indicar a ausência de valor: \texttt{null} e \texttt{undefined}. Embora eles sejam semelhantes em alguns aspectos, há diferenças importantes entre eles. O \texttt{null} é um valor atribuído explicitamente a uma variável para indicar a ausência intencional de qualquer objeto ou valor. Em outras palavras, quando uma variável é definida como \texttt{null}, ela indica que não há valor ou objeto associado a ela. Por outro lado, \texttt{undefined} é um valor atribuído automaticamente a uma variável que foi declarada, mas não recebeu um valor. Em outras palavras, uma variável é definida como \texttt{undefined} quando não foi inicializada com um valor válido.

O tipo \texttt{any} em \textit{TypeScript} é usado quando não se tem certeza sobre o tipo de uma variável ou quando a variável pode ter diferentes tipos ao longo do programa. Essencialmente, o tipo \texttt{any} desabilita o sistema de tipos do \textit{TypeScript} para uma variável específica, permitindo que ela aceite qualquer valor. Isso pode ser útil em situações em que você está migrando código \textit{JavaScript} existente para o \textit{TypeScript} e ainda não tem informações sobre o tipo adequado. No entanto, o uso excessivo do tipo \texttt{any} pode anular os benefícios de segurança de tipos que o \textit{TypeScript} oferece.

Por outro lado, antagônico ao tipo \texttt{any}, o tipo \texttt{never} em \textit{TypeScript} representa o conjunto de valores que nunca ocorre. Isso significa que uma função que retorna \texttt{never} nunca retorna ou sempre lança uma exceção. Esse tipo é útil em cenários onde você quer garantir que uma função nunca conclua com sucesso, como funções que lançam erros ou loops infinitos. No \textit{JavaScript}, não há um tipo diretamente equivalente a \texttt{never}. No entanto, você pode alcançar comportamento semelhante usando exceções ou loops infinitos.

\section{Números}
\label{JS_numeros}

Os números desempenham um papel fundamental na programação Web, pois são essenciais para executar cálculos, representar quantidades, armazenar informações e controlar o comportamento de aplicativos e sites. Através da manipulação de números, é possível realizar operações matemáticas complexas, como cálculos financeiros, estatísticos e de geometria, além de gerenciar datas e horários. Além disso, os números são cruciais para criar interações dinâmicas e responsivas, permitindo a validação de formulários, a geração de números aleatórios, a animação de elementos e muito mais. A habilidade de manipular números com precisão e eficiência é essencial para desenvolver aplicativos Web que sejam funcionais, confiáveis e capazes de fornecer uma experiência fluida aos usuários. Ambas as linguagens JavaScript (JS) e TypeScript (TS) fornecem recursos poderosos para manipular e trabalhar com números de maneira eficiente.

\textbf{Números Decimais:} JavaScript e \textit{TypeScript} suportam números decimais, representados pelo tipo \texttt{number}. Ambas as linguagens permitem realizar operações aritméticas com números decimais, como adição, subtração, multiplicação e divisão. O exemplo na \autoref{lst:decimals} mostra as operações de adição, subtração, multiplicação e divisão entre dois números \texttt{a} e \texttt{b}, respectivamente, na terceira, quarta, quinta e sexta linha.

\Needspace{0.3\textheight}
\begin{lstlisting}[language=TypeScript, caption={Operações com números decimais em \textit{JavaScript} e \textit{TypeScript}}, label={lst:decimals}]
let a = 5.0;
let b = 2.5;
let sum = a + b;     // Resultado: 7.5
let dif = a - b;     // Resultado: 2.5
let prod = a * b;    // Resultado: 12.5
let div = a / b;     // Resultado: 2.0
\end{lstlisting}

\textbf{Números Binários:} Enquanto \textit{JavaScript} não possui um literal específico para números binários, o \textit{TypeScript} introduziu a sintaxe de literais binários com o prefixo \texttt{0b}. Isso permite a representação direta de números binários, como \texttt{0b1010} para o número decimal 10, como é possível observar na primeira linha da \autoref{lst:bases-num}.

\textbf{Números Octais:} Assim como os números binários, o \textit{JavaScript} não possui um literal para números octais. No entanto, o \textit{TypeScript} adicionou suporte para literais octais com o prefixo \texttt{0o}. Por exemplo, \texttt{0o16} representa o número decimal 14, como mostra a segunda linha da \autoref{lst:bases-num}.

\textbf{Números Hexadecimais:} Ambas as linguagens \textit{JavaScript} e \textit{TypeScript} suportam números hexadecimais através do uso do prefixo \texttt{0x}. Na terceira linha da \autoref{lst:bases-num} é explanado um exemplo onde \texttt{0xFF} representa o número decimal 255.

\begin{lstlisting}[language=TypeScript, caption={Números binários em \textit{TypeScript}}, label={lst:bases-num}]
let numBinario = 0b1010;   // Numero decimal 10
let numOctal = 0o16;       // Numero decimal 14
let numHexadecimal = 0xFF; // Numero decimal 255
\end{lstlisting}

\textbf{Exponenciação/Notação Científica:} JavaScript e \textit{TypeScript} suportam a notação científica para representar números muito grandes ou muito pequenos. Por exemplo, \texttt{1.23e+6} representa o número 1230000 e \texttt{1.23e-6} representa o número 0,00000123, como mostra o exemplo na \autoref{lst:scientific}.

\begin{lstlisting}[language=TypeScript, caption={Notação científica em \textit{JavaScript} e \textit{TypeScript}}, label={lst:scientific}]
let bigNum = 1.23e+6;   // Numero 1230000
let smallNum = 1.23e-6; // Numero 0.00000123
\end{lstlisting}

\subsection{Objeto Number} 

\textit{JavaScript} fornece o objeto \texttt{Number} embutido, que oferece uma variedade de métodos úteis para trabalhar com números, que também é válido em \textit{TypeScript}. Por exemplo, o método \texttt{toFixed()} retorna uma representação formatada do número com uma quantidade específica de casas decimais. Na \autoref{lst:numberObject}, demonstra-se o uso do objeto \texttt{Number} para formatar o número 42 em uma string com duas casas decimais, resultando em ``\texttt{42.00}''.

\begin{lstlisting}[language=TypeScript, caption={Utilizando o objeto Number em \textit{JavaScript} e \textit{TypeScript}}, label={lst:numberObject}]
let num = 42;
let fixedDecimal = num.toFixed(2); 
// Numero formatado com 2 casas decimais: "42.00"
\end{lstlisting}

O uso do objeto \texttt{Number} permite realizar operações mais avançadas, como conversões de bases numéricas e manipulação precisa de valores numéricos. Ele oferece uma série de métodos e propriedades, como \texttt{toString()}, que retorna a string do número em certa base numérica passada como argumento, entre outros, que podem ser observados na \autoref{tab:numberObject}, que apresenta alguns dos principais métodos e atributos que podem ser utilizados com o objeto \texttt{Number}. Esses recursos fornecem funcionalidades adicionais para a manipulação e formatação de números em \textit{JavaScript} e \textit{TypeScript}, permitindo realizar operações mais avançadas e precisas~\cite{mdn_numbers_dates}.

\begin{table}
\centering
\caption{Métodos e Atributos do Objeto Number em JS e TS}
\label{tab:numberObject}
\begin{tabular}{|l|l|}
\hline
\textbf{Método/Atributo} & \textbf{Descrição} \\ \hline
toExponential() & String do número em notação exponencial. \\ \hline
toFixed(n) & String do número com n casas decimais. \\ \hline
toPrecision(n) & String do número com a precisão especificada n. \\ \hline
toString(n) & String do número na base n. \\ \hline
valueOf() & Valor numérico primitivo do objeto Number. \\ \hline
MAX\_VALUE & O maior número possível da linguagem. \\ \hline
MIN\_VALUE & O menor número possível da linguagem. \\ \hline
POSITIVE\_INFINITY & Representa o infinito positivo. \\ \hline
NEGATIVE\_INFINITY & Representa o infinito negativo. \\ \hline
NaN & Representa um valor que não é um número. \\ \hline
\end{tabular}
\end{table}



\subsection{NaN (Not a Number)}
Em \textit{JavaScript} e \textit{TypeScript}, o valor especial \texttt{NaN} (\textit{Not a Number}) é usado para representar um resultado numérico inválido ou indefinido. Ele é retornado quando ocorre uma operação matemática que não pode produzir um resultado numérico válido. O \texttt{NaN} é um valor do tipo número, mas não é igual a nenhum número real ou infinito. Ele é considerado um valor falso quando usado em um contexto booleano. Por exemplo, a expressão \texttt{NaN === NaN} retorna \texttt{false}, o que indica que \texttt{NaN} não é igual a ele mesmo. Alguns exemplos de casos em que o \texttt{NaN} pode ser gerado são:

\begin{itemize}
  \item Divisão por zero: Quando um número é dividido por zero, o resultado é indefinido e é representado como \texttt{NaN}.
  \item Operações matemáticas inválidas: Se uma operação matemática que não faz sentido é executada, como tentar obter a raiz quadrada de um número negativo, o resultado será \texttt{NaN}.
  \item Conversões numéricas inválidas: Ao tentar converter uma string que não representa um número válido em \textit{JavaScript} ou \textit{TypeScript}, o resultado será \texttt{NaN}.
\end{itemize}

O \texttt{NaN} tem algumas propriedades peculiares quando usado em cálculos, pois o \texttt{NaN} é contagioso. Qualquer operação aritmética com \texttt{NaN} resultará em \texttt{NaN}, como por exemplo: \texttt{NaN + 5} ou \texttt{NaN * 10}. Por isso, ao lidar com valores numéricos em \textit{JavaScript} e \textit{TypeScript}, é importante verificar se um valor é \texttt{NaN} antes de usá-lo. Pode-se usar a função \texttt{isNaN()} para verificar se um valor é ou não \texttt{NaN}. Essa função retorna \texttt{true} se o valor fornecido for \texttt{NaN} e \texttt{false} caso contrário.

Embora o comportamento do \texttt{NaN} seja semelhante em \textit{JavaScript} e \textit{TypeScript}, o \textit{TypeScript} oferece suporte adicional aos tipos de dados estáticos, o que pode ajudar a evitar erros relacionados ao uso de \texttt{NaN} em operações matemáticas. Com o \textit{TypeScript}, é possível definir tipos específicos para variáveis e parâmetros de função, o que pode ajudar a identificar erros de tipo em tempo de compilação e evitar situações em que \texttt{NaN} é usado incorretamente.




\section{Strings}
\label{JS_string}

As strings desempenham um papel crucial na programação Web, sendo fundamentais para manipular e exibir texto e conteúdo em aplicativos e sites. Elas permitem a representação e armazenamento de informações como nomes, mensagens, endereços e muitos outros dados textuais. Além disso, a manipulação de strings é essencial para realizar operações como concatenação, formatação, busca, substituição e extração de partes específicas do texto. Isso permite criar recursos como formulários interativos, validação de entrada de dados, filtragem e ordenação de listas, geração dinâmica de conteúdo e muito mais. A habilidade de trabalhar com strings de forma eficiente e precisa é essencial para o desenvolvimento de aplicativos Web robustos, funcionais e com uma interface de usuário envolvente e amigável. Ambas as linguagens \textit{JavaScript} e \textit{TypeScript} fornecem recursos poderosos para trabalhar com strings.

\textbf{Declaração de Strings:} Em \textit{JavaScript} e \textit{TypeScript}, as strings podem ser declaradas utilizando aspas simples (` ') ou aspas duplas (`` ''). Ambas as formas são igualmente válidas e produzem a mesma saída. Pode-se escolher usar aspas simples ou aspas duplas com base em preferências pessoais ou em convenções de estilo de código adotadas em um projeto específico. a \autoref{lst:dec-string} apresenta um exemplo onde se declara uma variável \texttt{mensagem} com o valor `Ola, mundo!', utilizando-se aspas simples, e outra variável \texttt{nome} com o valor ``Maria'', declarada com aspas duplas.

\begin{lstlisting}[language=TypeScript, label=lst:dec-string, caption=Declaração de String em JS e TS.]
let mensagem = 'Ola, mundo!';
let nome = "Maria";
\end{lstlisting}


\textbf{Concatenação de Strings:} A concatenação de strings é uma operação comum ao lidar com texto. Em \textit{JavaScript} e \textit{TypeScript}, pode-se concatenar strings utilizando o operador de adição (+) ou o método \texttt{concat()}. Para exemplificar, na \autoref{lst:concat-string}, declara-se uma variável \texttt{saudacao} com o valor ``Ola, '' e outras variáveis \texttt{nome1} e \texttt{nome2} com os valores ``Maria'' e ``Jose'', respectivamente. Em seguida, utiliza-se o operador de adição (+) para concatenar as strings \texttt{saudacao} e \texttt{nome1}, armazenando o resultado na variável \texttt{mensagem}, que resultou em ``Ola, Maria''. Para realizar outra concatenação, foi utilizado o método \texttt{concat()} entre as strings \texttt{saudacao} e \texttt{nome2}, resultando em ``Olá, Jose''.

\begin{lstlisting}[language=TypeScript, label=lst:concat-string, caption=Concatenção de String em JS e TS.]
let saudacao = "Ola, ";
let nome1 = "Maria";
let nome2 = "Jose";
let mensagem1 = saudacao + nome1;      // "Ola, Maria"
let mensagem2 = saudacao.concat(nome2); // "Ola, Jose"
\end{lstlisting}

\textbf{Acessando Caracteres:} Caracteres individuais de uma string podem ser acessados utilizando a notação de colchetes ([]). Os índices começam em zero, ou seja, o primeiro caractere tem o índice 0. No exemplo da \autoref{lst:acesso-caracteres}, uma variável \texttt{nome} é declarada com o valor ``Alice''. Em seguida, foi utilizada a notação de colchetes com o índice zero para acessar o primeiro caractere da string, que é ``A'', e ainda, foi utilizado o atributo \texttt{length} da string para pegar seu tamanho e posteriormente subtraído em uma unidade para indicar o índice do último caractere, que é ``e''.

\begin{lstlisting}[language=TypeScript, caption= Acessando caracteres de strings em JS e TS., label=lst:acesso-caracteres]
let nome = "Alice";
let primeiroCaractere = nome[0];             // "A"
let ultimoCaractere = nome[nome.length - 1]; // "e"
\end{lstlisting}

\textbf{Métodos de Manipulação de Strings:} \textit{JavaScript} e \textit{TypeScript} fornecem uma variedade de métodos embutidos para manipulação de strings. a \autoref{lst:manip-string} ilustra a utilização do método \texttt{toUpperCase} para transformar todos os caracteres da string ``\texttt{julia}'' em caracteres maiúsculos. Alguns dos métodos de manipulação de string em JS e TS mais comuns incluem:

\begin{itemize}
\item \texttt{length}: Retorna o comprimento da string.
\item \texttt{toUpperCase()}: Converte a string para letras maiúsculas.
\item \texttt{toLowerCase()}: Converte a string para letras minúsculas.
\item \texttt{substring(startIndex, endIndex)}: Retorna uma parte da string, começando do índice \texttt{startIndex} até o índice \texttt{endIndex} (não incluído).
\item \texttt{split(separator)}: Divide a string em um array de substrings com base no separador especificado.
\item \texttt{trim()}: Remove espaços em branco do início e do final da string.

\end{itemize}

\begin{lstlisting}[language=TypeScript, caption= Acessando caracteres de strings em JS e TS., label=lst:acesso-caracteres, label=lst:manip-string]
let nome = "julia";
let nomeMaiusculo = nome.toUpperCase() // "JULIA"
\end{lstlisting}


\section{Operações Básicas}
\label{JS_operacoesBasica}

As operações básicas são realizadas pelos operadores nativos da linguagem que são amplamente utilizados, são como comandos atômicos dentro da linguagem. Em \textit{JavaScript} (JS) e \textit{TypeScript} (TS), dentre essas operações estão: operadores aritméticos, operadores de atribuição e operadores de comparação.

\subsection{Atribuição}

Atribuição é uma operação que armazena um valor em uma variável. O operador utilizado para essa tarefa em JS e TS é o igual (=). Como mostra a \autoref{lst:atrib}, há a declaração de uma variável numérica e em seguida é-lhe atribuído o valor 5, na primeira e segunda linha, respectivamente. Também é possível fazer o mesmo processo de declaração e atribuição de uma forma mais sucinta, como é expresso na terceira linha.


\begin{lstlisting}[language=TypeScript, caption=Exemplo de Atribuição em JS e TS., label=lst:atrib]
let minhaVariavel;
minhaVariavel = 5;
let minhaOutraVariavel = 5;
\end{lstlisting}

A \emph{inicialização} é a primeira operação de atribuição que uma variável recebe, e é fundamental para garantir que a variável não fique com um valor vazio ou indefinido. No \textit{TypeScript}, ao inicializar uma variável com um valor, a linguagem realiza a \emph{inferência de tipo}, determinando automaticamente qual será o tipo da variável com base nessa atribuição. Se uma variável for declarada e receber um valor na mesma linha, o compilador é capaz de inferir seu tipo, dispensando a necessidade de declarar explicitamente o tipo da variável. No entanto, se uma variável for declarada sem especificar seu tipo e não receber um valor na mesma linha, seu tipo será assumido como \texttt{any} no código \textit{TypeScript}.

\subsection{Comentários}

Os comentários, no contexto da programação, são trechos do código os quais não interferem no programa, são apenas anotações no decorrer do código. Para iniciar um comentário, tanto em JS como em TS, basta inserir duas barras (//) no código e todo o texto que estiver em sequência na mesma linha será considerado como um comentário. Caso haja a necessidade de escrever um comentário que ocupe mais que uma linha, pode-se iniciá-lo com uma barra e um asterisco (/*) e finalizá-lo com um asterisco e uma barra (*/). Como mostra o exemplo na \autoref{lst:coment}, na primeira linha de código há um comando de declaração e inicialização de uma variável, e após esse comando, na mesma linha, um comentário ``Comentario em uma linha'', e após, nas linhas seguintes, um outro comentário ``Este comentario pode tomar varias linhas''.

\begin{lstlisting}[language=TypeScript, label=lst:coment, caption=Exemplo de Comentário em JS e TS.]
let minhaVariavel = 2; // Comentario em uma linha
/* Este comentario
   pode tomar
   varias linhas */
\end{lstlisting}

A importância dos comentários reside no fato de que eles facilitam a compreensão do código para os desenvolvedores, tanto para o próprio autor quanto para outros que possam trabalhar ou dar manutenção no código no futuro. Comentários bem escritos ajudam a descrever a lógica por trás de uma seção de código, a fornecer contextos importantes e a documentar decisões de design. Os comentários também são úteis para adicionar lembretes, fazer anotações ou desativar temporariamente partes do código durante o desenvolvimento. Além disso, eles podem ajudar na depuração, permitindo que os desenvolvedores expliquem o propósito de certas linhas de código ou identifiquem possíveis problemas. Ao utilizar comentários, é importante seguir algumas boas práticas, o que inclui: manter os comentários atualizados à medida que o código evolui, escrever comentários claros e concisos, evitar comentários óbvios ou redundantes e evitar comentários em excesso que possam poluir o código.

\subsection{Operações Aritméticas}

Algumas operações aritméticas básicas e recorrentemente utilizadas na programação são: adição, subtração, multiplicação, divisão e módulo (resto da divisão). Essas operações são realizadas em \textit{JavaScript} e \textit{TypeScript} de maneira semelhante à matemática convencional, e seus resultados podem ser armazenados em variáveis para posterior utilização. É importante mencionar que, assim como na matemática, os parênteses podem ser utilizados para definir a ordem de prioridade das operações.

a \autoref{lst:op-aritmeticas} mostra alguns exemplos de operações aritméticas básicas, dentre elas: na primeira linha, a adição dos valores 2 e 5, em que o resultado é atribuído à variável \texttt{x}, que terá o valor 7; na segunda linha, a subtração do valor 5 do valor 15, em que resultado é atribuído à variável \texttt{y}, que terá o valor 10; na terceira linha, a multiplicação dos valores 3 e 5, em que o resultado é atribuído à variável \texttt{z}, que terá o valor 15; na quarta linha, a divisão de 15 por 5, em que o resultado é atribuído à variável \texttt{w}, que terá o valor 3; e na quinta linha, a operação de resto da divisão, dividindo 7 por 5, em que o resultado é atribuído à variável \texttt{k}, que terá o valor 2.


\begin{lstlisting}[language=TypeScript, caption=Exemplos de Operações Aritméticas em \textit{JavaScript} e \textit{TypeScript}., label=lst:op-aritmeticas]
let x = 2 + 5; // x vale 7
let y = 15 - 5; // y vale 10
let z = 3 * 5; // z vale 15
let w = 15 / 5; // w vale 3
let k = 7 % 5; // k vale 2
\end{lstlisting}

Essas operações aritméticas são úteis em uma variedade de cenários de programação, desde cálculos simples até operações mais complexas envolvendo variáveis e expressões. É importante ter em mente a ordem de precedência dos operadores aritméticos, bem como o uso adequado de parênteses quando necessário, para obter os resultados desejados.


%%TODO
\subsection{Objeto Math}

O objeto \texttt{Math} em \textit{JavaScript} e \textit{TypeScript} fornece um conjunto de propriedades e métodos para realizar operações matemáticas mais avançadas. Essas operações abrangem desde funções trigonométricas e exponenciais até operações de arredondamento e geração de números aleatórios. Nessa subseção são explanados alguns dos principais métodos e propriedades disponíveis no objeto \texttt{Math}.

\textbf{Métodos Trigonométricos:} O objeto \texttt{Math} oferece métodos para calcular funções trigonométricas como seno, cosseno e tangente. Na \autoref{lst:math-sin}, usa-se o método \texttt{Math.sin()} para calcular o seno de um ângulo. No exemplo, o ângulo é definido como $\frac{\pi}{4}$, que é equivalente a 45 graus em radianos. O resultado é armazenado na variável \texttt{sineValue} e, em seguida, é impresso no console. Além do \texttt{Math.sin()}, o objeto \texttt{Math} também oferece os métodos \texttt{Math.cos()} e \texttt{Math.tan()} para calcular o cosseno e a tangente de um ângulo, respectivamente.

\Needspace{0.3\textheight}
\begin{lstlisting}[language=TypeScript, caption=Exemplo de uso do método \texttt{Math.sin()}., label=lst:math-sin]
// Angulo de 45 graus em radianos
const angle = Math.PI / 4; 
// Calcula o seno do angulo
const sineValue = Math.sin(angle); 
\end{lstlisting}



\textbf{Métodos Exponenciais:} O objeto \texttt{Math} possui métodos para cálculos exponenciais, como calcular a potência de um número e calcular o logaritmo. Na \autoref{lst:math-pow}, o método \texttt{Math.pow()} é utilizado para calcular a potência de um número. No exemplo, eleva-se o número 2 à potência 3, o que resulta em 8. O resultado é armazenado na variável \texttt{result} e é impresso no console. Além do \texttt{Math.pow()}, o objeto \texttt{Math} também possui o método \texttt{Math.sqrt()} para calcular a raiz quadrada de um número e o método \texttt{Math.log()} para calcular o logaritmo natural.

\Needspace{0.3\textheight}
\begin{lstlisting}[language=TypeScript, caption=Exemplo de uso do método \texttt{Math.pow()}., label=lst:math-pow]
const base = 2;
const exponent = 3;
// Calcula 2 elevado a 3, que resulta em 8
const result = Math.pow(base, exponent); 
\end{lstlisting}



\textbf{Métodos de Arredondamento:} O objeto \texttt{Math} fornece métodos para arredondar números para cima, para baixo e para o valor inteiro mais próximo. Na \autoref{lst:math-round}, utiliza-se os métodos \texttt{Math.ceil()}, \texttt{Math.floor()} e \texttt{Math.round()} para arredondar um número. No exemplo, o número inicial é 4.8. O método \texttt{Math.ceil()} é usado para arredondar para cima, \texttt{Math.floor()} para arredondar para baixo e \texttt{Math.round()} para arredondar para o valor inteiro mais próximo. Os resultados são armazenados em variáveis diferentes e são impressos no console.

\begin{lstlisting}[language=TypeScript, caption=Exemplo de uso dos métodos de arredondamento em \texttt{Math}., label=lst:math-round]
const number = 4.8;
// Arredonda para cima: 5
const roundedUp = Math.ceil(number); 
// Arredonda para baixo: 4
const roundedDown = Math.floor(number); 
// Arredonda para o valor inteiro mais proximo: 5
const rounded = Math.round(number); 
\end{lstlisting}



\textbf{Métodos de Números Aleatórios:} O objeto \texttt{Math} também possui um método para gerar números aleatórios. Na \autoref{lst:math-random}, foi utilizado o método \texttt{Math.random()} para gerar um número aleatório entre 0 e 1. O resultado é armazenado na variável \texttt{random} e é impresso no console.

\begin{lstlisting}[language=TypeScript, caption=Exemplo de uso do método \texttt{Math.random()}., label=lst:math-random]
// Gera um numero aleatorio entre 0 e 1
const random = Math.random(); 
\end{lstlisting}

%% https://developer.mozilla.org/pt-BR/docs/Web/JavaScript/Guide/Numbers_and_dates


\subsection{Operações Relacionais}

As operações relacionais são utilizadas para realizar comparações entre valores. Essas operações são comumente usadas em estruturas de condição, como os comandos \texttt{if} e \texttt{switch} (que serão apresentados posteriormente), para verificar se uma determinada condição é verdadeira ou falsa. Existem seis operações relacionais em \textit{JavaScript} e \textit{TypeScript}: igual (==), diferente (\texttt{!=}), maior que (\texttt{>}), maior ou igual (\texttt{>=}), menor que (\texttt{<}) e menor ou igual (\texttt{<=}). Quando se utiliza uma dessas operações para comparar dois valores, o resultado será um valor booleano (verdadeiro/\textit{true} ou falso/\textit{false}).

Na \autoref{lst:op-relacionais} há alguns exemplos de uso das operações relacionais entre as variáveis \texttt{a} e \texttt{b}, dentre elas: na terceira linha, a operação de igualdade (==), que resultará em \textit{false}, que é atribuído à variável \texttt{x}; na quarta linha, a operação de diferença, o resultado será atribuído à variável \texttt{y}, que terá o valor booleano \textit{true}; na quinta linha, a operação de maior que, que resultará em \textit{true} e será atribuído à variável \texttt{z}; na sexta linha, a operação de maior ou igual, que resultará em \textit{true} e será atribuído à variável \texttt{w}; na sétima linha, a operação de menor que, que resultará em \textit{false} e será atribuído à variável \texttt{k}; e por fim, na oitava linha, a operação de menor ou igual, atribuída à variável \texttt{m}, que valerá \textit{false}.

\begin{lstlisting}[language=TypeScript, caption=Exemplo de Operações Relacionais em \textit{JavaScript} e \textit{TypeScript}., label=lst:op-relacionais]
const a = 5;
const b = 4;
const x = a == b; // x = false 
const y = a != b; // y = true
const z = a > b;  // z = true
const w = a >= b; // w = true
const k = a < b;  // k = false
const m = a <= b; // m = false
\end{lstlisting}

Além desses operadores, também existem os operadores estritamente igual, os quais são definidos sintaticamente como três iguais juntos (\texttt{===}),  e o estritamente diferente, com sintaxe de um ponto de exclamação seguido de dois sinais de igual (\texttt{!==}). Esses operadores diferem dos operadores de igualdade (\texttt{==}) e diferença (\texttt{!=}) por considerarem também os tipos dos valores sendo comparados. Eles exigem que os valores sejam iguais em valor e tipo para serem considerados verdadeiros. Esses operadores são úteis quando se deseja garantir que os valores sendo comparados sejam idênticos em valor e tipo. 

É recomendado utilizar os operadores estritamente igual e estritamente diferente sempre que possível, pois eles oferecem uma comparação mais precisa e ajudam a evitar resultados inesperados devido à coerção de tipo. Além disso, é importante notar que, ao realizar comparações em \textit{JavaScript} e \textit{TypeScript}, é necessário levar em consideração as regras de coerção de tipo (\textit{type coercion}) que podem ocorrer durante as operações. Por exemplo, ao comparar um número com uma string, o \textit{JavaScript} tentará converter a string em um número antes de realizar a comparação.


\subsection{Entrada e Saída}

Os comandos básicos de entrada e saída são utilizados para interagir com o usuário ou com dispositivos externos, permitindo que o programa receba informações e exiba resultados. Em \textit{JavaScript} e \textit{TypeScript} os comandos de entrada e saída são os mesmos. Para fazer a leitura de dados do usuário, utiliza-se o comando \texttt{prompt}, que exibe uma janela de diálogo na qual o usuário pode inserir um valor. Já para exibir informações para o usuário, se utiliza o comando \texttt{console.log}, que imprime uma mensagem no console do navegador ou no console do \textit{Node.js}. a \autoref{lst:input-output} ilustra a utilização dos comandos, em que na segunda linha é solicitado o nome do usuário com a função \texttt{prompt} e depois imprime no console o nome digitado com o comando \texttt{console.log}.

\begin{lstlisting}[language=TypeScript, caption=Comando \texttt{prompt} e \texttt{console.log} em JS e TS., label=lst:input-output]
// solicita uma entrada do usuario
let nome = prompt("Digite seu nome:");
// imprime no console a entrada do usuario
console.log(nome);   
\end{lstlisting}

Além do comando \texttt{console.log} que imprime mensagens no console, existe também o comando \texttt{alert}, que exibe uma janela de alerta no navegador com uma mensagem para o usuário, como mostrado na \autoref{fig:ex-exec-js}. O \texttt{alert} é um comando de saída muito utilizado para exibir informações importantes ou notificações para o usuário. Esse comando é particularmente útil em contextos onde é necessário chamar a atenção do usuário para alguma informação importante, como validações de formulários, confirmações de ações ou exibição de erros. No entanto, é importante usar o \texttt{alert} com moderação, pois janelas de alerta excessivas podem ser intrusivas e perturbar a experiência do usuário.

Em \textit{JavaScript} e \textit{TypeScript}, o uso do \texttt{alert} é semelhante ao \texttt{console.log}. a \autoref{lst:alert-output}, baseado na \autoref{lst:input-output}, ilustra o uso do \texttt{alert} para exibir o nome digitado pelo usuário em uma janela de alerta. Nesse exemplo, após solicitar o nome do usuário com o comando \texttt{prompt}, o comando \texttt{alert} é utilizado para exibir o nome em uma janela de alerta. Assim, o usuário verá o valor digitado no alerta do navegador.

\begin{lstlisting}[language=TypeScript, caption=Exemplo de Utilização do Comando Alert em JS e TS., label=lst:alert-output]
// solicita uma entrada do usuario
let nome = prompt("Digite seu nome:");
// exibe o nome em uma janela de alerta
alert(nome);
\end{lstlisting}


\section{Condicionais}
\label{JS_condicionais}

As condicionais são estruturas de controle que permitem que o programa execute diferentes blocos de código com base em condições específicas. Elas são amplamente utilizadas em \textit{JavaScript} e \textit{TypeScript} para controlar o fluxo do programa com base em valores e expressões. As condicionais mais comuns nessas linguagens são \texttt{if}/\texttt{else}, mas também existe a estrutura \texttt{switch}/\texttt{case} que possibilita, também, realizar o controle de fluxo do código para diferentes valores de uma variável.

\subsection{\texttt{if}/\texttt{else}}

O comando \texttt{if} é usado para testar uma condição e executar um bloco de código caso a condição seja verdadeira. a \autoref{lst:if} exemplifica a utilização da estrutura, onde o programa verifica se a variável \texttt{idade} é maior ou igual a 18. Se essa condição for verdadeira, o bloco de código dentro do \texttt{if} é executado e exibe a mensagem ``Voce eh maior de idade'' no console.

\begin{lstlisting}[language=TypeScript, caption=Exemplo de condicional \texttt{if} em JS e TS., label=lst:if]
if (idade >= 18) {
    console.log("Voce eh maior de idade");
}
\end{lstlisting}

O comando \texttt{else} é usado em conjunto com o \texttt{if} e permite executar um bloco de código caso a condição no \texttt{if} seja falsa. Como é possível observar na \autoref{lst:ifelse}, se a condição no \texttt{if} for falsa, o bloco de código dentro do \texttt{else} é executado e exibe a mensagem ``Voce eh menor de idade'' no console.


\begin{lstlisting}[language=TypeScript, caption=Exemplo de condicional \texttt{if/else} em JS e TS., label=lst:ifelse]
if (idade >= 18) {
    console.log("Voce eh maior de idade");
} else {
    console.log("Voce eh menor de idade");
}
\end{lstlisting}

A estrutura \texttt{else if} é usada para testar condições adicionais, caso a condição no \texttt{if} seja falsa. Na \autoref{lst:ifelseif}, se a condição no \texttt{if} for falsa, o programa testa a condição no \texttt{else if} e, se essa condição for verdadeira, o bloco de código correspondente é executado e exibe a mensagem ``Voce tem idade suficiente para dirigir'', caso contrário, o bloco de código no \texttt{else} é executado e exibe a mensagem ``Voce eh menor de idade''.

\begin{lstlisting}[language=TypeScript, caption=Exemplo de condicional utlizando \texttt{if/else} e \texttt{else if} em JS e TS., label=lst:ifelseif]
if (idade >= 18) {
    console.log("Voce eh maior de idade.");
} else if (idade >= 16) {
    console.log("Voce tem idade suficiente para dirigir.");
} else {
    console.log("Voce eh menor de idade.");
}
\end{lstlisting}

\subsection{\texttt{switch}/\texttt{case}}

A estrutura \texttt{switch/case} é útil quando há várias condições a serem testadas e se quer evitar usar várias estruturas \texttt{if/else}. O comando \texttt{switch} inicia a estrutura, seguido pela variável que será testada nos diferentes casos. Como é possível observar na \autoref{lst:switchcase}, o programa testa a variável \texttt{diaDaSemana} e executa o bloco de código correspondente a cada caso. Se o valor for ``Sabado'' ou ``Domingo'', a mensagem ``Hoje eh fim de semana'' será exibida no console. Caso contrário, o bloco de código no caso \texttt{default} é executado e exibe a mensagem ``Hoje eh dia de trabalho''.

\begin{lstlisting}[language=TypeScript, caption=Exemplo de condicional \texttt{switch/case} em JS e TS., label=lst:switchcase]
switch (diaDaSemana) {
    case "Sabado":
    case "Domingo":
        console.log("Hoje eh fim de semana");
        break;
    default:
        console.log("Hoje eh dia de trabalho");
        break;
}
\end{lstlisting}

O comando \texttt{break} é uma instrução utilizada em JavaScript e TypeScript para interromper a execução de um laço de repetição ou de um bloco de código dentro de uma estrutura condicional. Quando o comando \texttt{break} é executado, o controle do programa é transferido imediatamente para fora do laço ou bloco de código, ignorando as instruções restantes. É importante notar que cada bloco de código dentro de \texttt{case} deve terminar com o comando \texttt{break}. Isso impede que o programa execute os casos seguintes, mesmo que a condição seja verdadeira.


\section{Laços de Repetição}
\label{JS_lacosRepeticao}

Os laços de repetição são estruturas que permitem que o programa execute um bloco de código várias vezes, com base em uma condição específica. Eles são amplamente utilizados em \textit{JavaScript} e \textit{TypeScript} para automatizar tarefas repetitivas e percorrer listas de elementos. Existem dois tipos principais de laços de repetição nessas linguagens: \texttt{for} e \texttt{while}. Cada um tem sua utilidade, dependendo da situação.

\subsection{For}

O laço \texttt{for} é frequentemente usado quando se conhece o número exato de vezes que se deseja repetir o bloco de código. Ele possui três partes: a inicialização, a condição e o incremento. A inicialização é onde a variável do laço é definida e inicializada. A condição é a expressão que deve ser verdadeira para continuar a execução do bloco de código. O incremento define como a variável do laço é atualizada a cada iteração.

\begin{lstlisting}[language=TypeScript, caption=Exemplo do laço de repetição \texttt{for} em JS e TS., label=lst:for]
for (let i = 0; i < 10; i++) {
    console.log(i);
}
\end{lstlisting}

Na \autoref{lst:for}, o programa executa o bloco de código dentro do laço de repetição dez vezes. Ele exibe os números de 0 a 9 no console, pois o que acontece é o seguinte: na primeira repetição, o valor de \texttt{i} assume o valor da sua inicialização ``\texttt{let i = 0}'', e executa o bloco de código dentro do loop; após finalizar a execução do bloco, o incremento ``\texttt{i++}'' é realizado, e \texttt{i} recebe seu valor anterior somado de um; depois do incremente a verificação da condição ``\texttt{i < 10}'' é realizada e, se for verdadeira, o bloco de código do \texttt{for} é repetido, caso não, o \textit{loop} termina. Esse processo de incremento e verificação de condição é realizado em toda iteração.


\subsection{While}

O laço \texttt{while} é comumente usado quando não se sabe antecipadamente quantas vezes é necessário repetir o bloco de código. Ele repete o bloco de código enquanto a condição especificada for verdadeira. Como mostra a \autoref{lst:while}, o programa executa o bloco de código dentro do laço de repetição enquanto a variável \texttt{i} for menor que 10. Ele exibe os números de 0 a 9 no console, incrementando a variável \texttt{i} a cada iteração.

\Needspace{0.3\textheight}
\begin{lstlisting}[language=TypeScript, caption=Exemplo do laço de repetição \texttt{while} em JS e TS., label=lst:while]
let i = 0;
while (i < 10) {
    console.log(i);
    i++;
}
\end{lstlisting}

É importante garantir que a condição do laço de repetição eventualmente se torne falsa, caso contrário, o laço será executado infinitamente, causando um loop infinito no programa. Portanto, é necessário cuidado ao definir a condição para evitar esse problema.





\section{Funções}
\label{JS_funcoes}

Funções são blocos de código que podem ser definidos em um programa para realizar uma tarefa específica. Em \textit{TypeScript}, as funções podem ser definidas e usadas de maneira semelhante ao \textit{JavaScript}, mas com a adição de tipos de dados estáticos. Isso torna as funções em \textit{TypeScript} mais robustas em termos de verificação de tipos e fornecem uma maneira mais clara de entender os tipos de entrada e saída esperados.

\textbf{Declaração e Execução de Funções:} Para se declarar uma função usa-se a palavra-chave \texttt{function}, seguida do nome da função, uma lista de parâmetros entre parênteses e o corpo da função dentre chaves, como mostra a \autoref{lst:function-js}. No caso do \textit{TypeScript} pode-se definir ainda o tipo de dado do retorno da função, como elucida o exemplo da \autoref{lst:function-ts}. Já para executar uma função declarada, basta chamar a função pelo seu nome e passar os devidos parâmetros, caso necessário, como mostra a quarta linha de ambos os exemplos em JS e TS.

\begin{lstlisting}[language=TypeScript, caption=Exemplo de Declaração de Função em JS., label=lst:function-js]
function saudacao() {
  console.log("Ola, mundo!");
}
saudacao();  // Log: "Ola, mundo!"
\end{lstlisting}

\begin{lstlisting}[language=TypeScript, caption=Exemplo de Declaração de Função em TS., label=lst:function-ts]
function saudacao(): void {
  console.log("Ola, mundo!");
}
saudacao();  // Log: "Ola, mundo!"
\end{lstlisting}



\textbf{Parâmetros da Função:} Ao declarar uma função, pode-se especificar seus parâmetros. Na linguagem \textit{TypeScript} é possível ainda especificar os tipos dos parâmetros, ajudando a melhorar a verificação de tipos. No exemplos da \autoref{lst:ex-func-js} e \autoref{lst:ex-func-ts}, pode-se observar a especificação de duas variáveis parâmetro \texttt{a} e \texttt{b}, que pelo contexto, espera-se que sejam do tipo numérico. Ou seja, ao chamar a função \texttt{somar}, deve-se passar dois valores numéricos para que a função use esses valores para executar.

\begin{lstlisting}[language=TypeScript, caption=Exemplo de Declaração de Função com Parâmetros em JS., label=lst:ex-func-js]
function somar(a, b): void {
  console.log(a + b);
}
\end{lstlisting}

\begin{lstlisting}[language=TypeScript, caption=Exemplo de Declaração de Função com Parâmetros em TS., label=lst:ex-func-ts]
function somar(a: number, b: number): void {
  console.log(a + b);
}
\end{lstlisting}

O exemplo da \autoref{lst:exec-ex-func} mostra alguns casos de execução das funções declaradas na \autoref{lst:ex-func-js} e \autoref{lst:ex-func-ts} em JS e TS. Pode-se observar a execução correta das somas, de 3 mais 5 que resulta em 8 e da adição de -10 a 10, que resulta em zero. Contudo, as diferenças na utilização das linguagens \textit{JavaScript} e \textit{TypeScript} começam a surgir mais fortemente a partir deste ponto, pois, caso sejam passados parâmetros de tipos que não sejam numéricos para a função \texttt{somar}, o verificador de tipos do \textit{TypeScript} alertará sobre a inconsistência no código, o que não ocorre em \textit{JavaScript}.



\Needspace{0.3\textheight}
\begin{lstlisting}[language=TypeScript, caption=Exemplos de Execução de Função em JS e TS., label=lst:exec-ex-func]
somar(3, 5);       // Resultado em JS e TS: 8
somar(10, -10);    // Resultado em JS e TS: 0
somar("JS", "TS"); // Resultado em JS: "JSTS"
/* Resultado em TS:  Error: Argument of type 'string' 
is not assignable to parameter of type 'number'. */
\end{lstlisting}



\textbf{Parâmetros Opcionais:} Em \textit{TypeScript} é possível também declarar uma função com parâmetro opcional, o que significa que ele pode ser passado ou não para a função. Para isso, usa-se o operador interrogação (\texttt{?}) após o nome do parâmetro, como evidencia o exemplo na \autoref{lst:opt-param}. Na nona linha são passados os três parâmetros da função e a soma de 3 mais 1 mais 4 ocorre e é mostrada o resultado 8 no console. Já na décima linha somente dois parâmetros, 5 e 2, são passados para a função \texttt{somar}, e a função executa sem erros mostrando o valor 7 no console, pois o terceiro parâmetro é opcional.


\begin{lstlisting}[language=TypeScript, caption=Exemplos de Parâmetros Opcionais em TS., label=lst:opt-param]
function somar(a:number, b:number, c?:number):void{
    if (c){
        console.log(a+b+c);
    }
    else{
        console.log(a+b);
    }
}
somar(3, 1, 4); // Log: 8
somar(5, 2);    // Log: 7
\end{lstlisting}

\textbf{Valores Padrão de Parâmetros:} Ainda na declaração das funções, é possível definir valores padrões para os parâmetros das funções nas linguagens JS e TS. Os parâmetros assumem os valores padrão caso, na chamada de execução da função, eles não sejam especificados. Isso viabiliza uma maneira de tornar os parâmetros opcionais. a \autoref{lst:default-params} ilustra essa ideia, onde na quarta linha os dois parâmetros são informados, 1 e 2, resultando em 3. Diferentemente, na quinta linha somente um dos parâmetros é informado e a função executa normalmente alocando o valor zero (valor padrão) ao segundo parâmetro. Já na sexta linha nenhum parâmetro é informado e todos os parâmetros, por padrão, recebem o valor zero, fazendo a soma resultar em zero.

\Needspace{0.3\textheight}
\begin{lstlisting}[language=TypeScript, caption=Exemplos de Valores Padrão de Parâmetros em JS., label=lst:default-params]
function somar(a = 0, b = 0){
    console.log(a+b+c);
}
somar(1, 2);   // Log: 3
somar(1);      // Log: 1
somar();       // Log: 0
\end{lstlisting}

\textbf{Retornos da Função:} O retorno de uma função é o valor que a finalização de sua execução assume. Em \textit{TypeScript} é possível especificar o tipo de retorno de uma função usando a notação de dois pontos (:) seguido do tipo de dado desejado. Se a função não retornar nenhum valor, pode-se usar o tipo void, como mostrado nos exemplos anteriores. Tomando como base ainda o exemplo da \autoref{lst:default-params}, que ao rodar a função \texttt{somar} com diferentes parâmetros, o resultado da soma dos parâmetros é mostrado na tela. Contudo, a função não tem um retorno, o que pode ser evidenciado no exemplo da \autoref{lst:func-return}. Nota-se então que a função \texttt{somar} mostrou o resultado, mas esse valor só existe dentro da execução da função e não pode ser resgatado externamente a ela. O retorno desempenha esse papel com o comando \texttt{return}, como mostra o exemplo da \autoref{lst:ex-return-func}.

\Needspace{0.3\textheight}
\begin{lstlisting}[language=TypeScript, caption=Evidenciando Retorno da Função., label=lst:func-return]
const a = somar(3, 5); // Log: 8
console.log(a);        // Log: undefined
\end{lstlisting}

\begin{lstlisting}[language=TypeScript, caption=Exemplo de Retorno de Função., label=lst:ex-return-func]
function somar(a = 0, b = 0){
    return a+b;
}
somar(3, 5); // nada aparece na tela
const x = somar(3, 5);
console.log(x); // Log: 8
\end{lstlisting}

\subsection{Arrow Function}

As \textit{arrow functions}, ou funções de seta, tem uma sintaxe mais curta e concisa para criar funções em JavaScript e TypeScript. Ao invés de escrever a palavra 'function' e o bloco de código da função, pode-se simplesmente usar a seta  \verb|(=>)| após os parâmetros da função para definir o bloco de código. Essa sintaxe simplificada traz benefícios em termos de legibilidade e escrita de código mais conciso. Como apresentado na \autoref{lst:arrow-func}, há uma função de soma declarada como uma \textit{arrow function}. A função recebe dois parâmetros opcionais, 'a' e 'b', e exibe a soma dos dois valores no console. O uso da \textit{arrow function} torna a declaração da função mais compacta e clara, permitindo que o código seja lido e compreendido de forma mais eficiente.

\begin{lstlisting}[language=TypeScript, caption=Exemplo de Arrow Function., label=lst:arrow-func]
// Arrow function
const soma = (a = 0, b = 0) => {
    console.log(a + b);
}    
soma(1, 2); // Log: 3
\end{lstlisting}

Declarar arrow functions como \texttt{const} é uma prática recomendada. Ao atribuir uma arrow function a uma constante, garante-se que a referência à função não será reatribuída a outro valor. Isso ajuda a evitar erros acidentais, pois impede que a função seja inadvertidamente substituída por outro valor. Além disso, utilizar \texttt{const} para funções arrow ajuda a transmitir a intenção de que a função é imutável e não deve ser modificada posteriormente.



\section{Tratamento de Exceções}
\label{JS_tratamentoExcecoes}

O tratamento de exceções desempenha um papel crucial na robustez e na confiabilidade do software. Ele evita que erros não tratados interrompam abruptamente a execução do programa, garantindo que o código seja capaz de lidar com falhas e tomar ações apropriadas. Além disso, ao capturar e tratar exceções, é possível fornecer mensagens de erro significativas e informativas, facilitando a depuração e a identificação de problemas. O tratamento de exceções pode ser realizado com os blocos \texttt{try}, \texttt{catch} e \texttt{continue} em linguagens de programação como \textit{JavaScript} e \textit{TypeScript}. Essa abordagem permite capturar exceções, tratar erros de forma controlada e continuar a execução do programa mesmo após um erro. 

O bloco \texttt{try} é usado para envolver o código que pode gerar exceções. Se uma exceção ocorrer dentro desse bloco, o fluxo de execução é imediatamente desviado para o bloco \texttt{catch}. Isso permite que você capture a exceção e forneça um tratamento adequado. O bloco \texttt{catch} contém o código que será executado quando uma exceção é lançada no bloco \texttt{try}. Nele pode-se escrever lógica personalizada para lidar com o erro. É possível fornecer um identificador de exceção após a palavra-chave \texttt{catch}, como \texttt{catch(erro)}, para acessar informações adicionais sobre a exceção, como a mensagem de erro ou detalhes específicos. A instrução \texttt{continue} é usada para pular para a próxima iteração de um loop, ignorando temporariamente o erro. Ela é útil quando se deseja continuar a execução do programa mesmo após um erro, sem interromper completamente o fluxo.


Na \autoref{lst:try-catch-continue-example}, tem-se um array de números e um valor de destino que se quer encontrar. No laço de repetição \texttt{for}, utiliza-se o bloco \texttt{try} para processar cada número do array. Se o valor de destino for encontrado, uma mensagem é exibida e usa-se \texttt{continue} para pular para a próxima iteração. Caso ocorra algum erro durante o processamento de um número, o bloco \texttt{catch} é acionado e exibe uma mensagem de erro.




\begin{lstlisting}[language=TypeScript, caption={Exemplo de uso de \texttt{try}, \texttt{catch} e \texttt{continue}} em JS e TS, label=lst:try-catch-continue-example]
const numbers = [1, 2, 3, 4, 5];
const targetValue = 7;

for (let i = 0; i < numbers.length; i++) {
  try {
    if (numbers[i] === targetValue) {
      console.log("Valor encontrado!");
      continue;
    }
    console.log("Processando numero: " + numbers[i]);
  } catch (erro) {
    console.log("Ocorreu um erro: " + erro.message);
  }
}
\end{lstlisting}

Já é sabido que o \textit{TypeScript} possui uma funcionalidade adicional em relação ao \textit{JavaScript}: a verificação de tipos estáticos. Isso significa que, ao utilizar o \texttt{try}, \texttt{catch} e \texttt{continue} no \textit{TypeScript}, pode-se fornecer tipos específicos para as exceções capturadas e também para as variáveis usadas dentro dos blocos \texttt{try} e \texttt{catch}, garantindo uma maior segurança e clareza no código.



\section{\textit{Promises}}
\label{JS_promises}

As \textit{Promises} são um recurso poderoso para lidar com operações assíncronas em \textit{JavaScript} e \textit{TypeScript}. Elas permitem executar tarefas assíncronas e lidar com seus resultados ou erros de forma mais estruturada. As \textit{Promises} foram introduzidas na especificação ES6 (ECMAScript 2015) para tratar assincronicidade de forma mais eficiente em \textit{JavaScript}. Uma \textit{Promise} representa um valor que pode estar disponível imediatamente, no futuro ou nunca. Ela pode estar em um dos três estados: \textit{pending} (pendente), \textit{fulfilled} (cumprida) ou \textit{rejected} (rejeitada).

Na \autoref{lst:promise-js}, uma nova \textit{Promise} é criada com uma função que recebe dois parâmetros: \texttt{resolve} e \texttt{reject}. Dentro dessa função, pode-se realizar operações assíncronas, como chamadas de API ou acesso a um banco de dados. Ao concluir a operação com sucesso, chama-se \texttt{resolve} com o valor resultante. Se ocorrer um erro, chama-se \texttt{reject} com o erro correspondente.

\Needspace{0.3\textheight}
\begin{lstlisting}[language=TypeScript, caption={Criação de uma Promise em JavaScript}, label=lst:promise-js]
const promise = new Promise((resolve, reject) => {
    // Logica assincrona aqui
    // Se a operacao for bem-sucedida, chame resolve(valor)
    // Se a operacao falhar, chame reject(erro)
});
\end{lstlisting}


O \textit{TypeScript}, sendo um superconjunto do \textit{JavaScript}, oferece suporte completo às \textit{Promises}. Além das funcionalidades disponíveis em \textit{JavaScript}, o \textit{TypeScript} adiciona a capacidade de especificar tipos para os valores resolvidos e rejeitados. Ao utilizar \textit{Promises} em \textit{TypeScript}, pode-se definir o tipo dos valores de retorno, proporcionando uma maior segurança e clareza no código. Na \autoref{lst:promise-ts}, a função \texttt{fetchData} retorna uma \textit{Promise} que deve resolver um valor do tipo \texttt{string}. Essa especificação de tipo permite ao \textit{TypeScript} fornecer verificação de tipo adicional e auxiliar os desenvolvedores durante o desenvolvimento.
\Needspace{0.3\textheight}

\begin{lstlisting}[language=TypeScript, caption={Criação de uma \textit{Promise} em \textit{TypeScript}}, label=lst:promise-ts]
function fetchData(): Promise<string> {
    return new Promise((resolve, reject) => {
        // Logica assincrona aqui
        // Se a operacao for bem-sucedida, chame resolve(valor)
        // Se a operacao falhar, chame reject(erro)
    });
}
\end{lstlisting}


Para encadear operações assíncronas com \textit{Promises} em \textit{JavaScript}, os métodos \texttt{then} e \texttt{catch} são úteis. O método \texttt{then} é chamado quando a \textit{Promise} é cumprida, recebendo o valor resolvido como argumento. O método \texttt{catch} é chamado quando a \textit{Promise} é rejeitada, recebendo o erro como argumento. Da mesma forma que em \textit{JavaScript}, pode-se encadear operações assíncronas com \textit{Promises} em \textit{TypeScript} usando os métodos \texttt{then} e \texttt{catch}. A diferença é que o \textit{TypeScript} pode inferir e aplicar os tipos corretos automaticamente, tornando o código mais legível e evitando erros de tipo.

É possível observar um exemplo na \autoref{lst:promise-usage}, onde a função \texttt{fetchData} retorna uma \textit{Promise} que resolve uma string simulando a obtenção de dados assíncronos. Em seguida, o método \texttt{then} é encadeado para lidar com o valor resolvido e o método \texttt{catch} para tratar qualquer erro que possa ocorrer durante a execução da \textit{Promise}. O \textit{TypeScript} garante que o valor retornado pela \textit{Promise} seja do tipo especificado (\texttt{string} neste caso) e também fornece verificação de tipo nos argumentos das funções passadas para \texttt{then} e \texttt{catch}, evitando erros de tipo durante a compilação.

\Needspace{0.3\textheight}
\begin{lstlisting}[language=TypeScript, caption={Exemplo de uso de \textit{Promises} em \textit{TypeScript}}, label=lst:promise-usage]
function fetchData(): Promise<string> {
    return new Promise((resolve, reject) => {
        // Simulando uma operacao assincrona
        setTimeout(() => {
            const data = 'Dados obtidos com sucesso!';
            resolve(data);
        }, 2000);
    });
}

fetchData()
.then((data) => {
    console.log(data);
})
.catch((error) => {
    console.error(error);
});
\end{lstlisting}

As \textit{Promises} são uma ferramenta essencial para lidar com operações assíncronas em \textit{JavaScript} e \textit{TypeScript}. O \textit{TypeScript}, com sua tipagem estática e recursos avançados de inferência de tipos, oferece uma experiência de desenvolvimento ainda mais robusta ao trabalhar com \textit{Promises}, permitindo detectar erros de tipo em tempo de compilação e fornecendo uma melhor documentação e compreensão do código.



\section{\textit{Callback}}
\label{JS_callback}

\textit{Callbacks} são um padrão de programação muito comum em \textit{JavaScript} e \textit{TypeScript} para lidar com operações assíncronas e chamadas de retorno. Uma função de \textit{callback} é uma função que é passada como argumento para outra função e é executada posteriormente, geralmente quando uma operação assíncrona é concluída ou um evento ocorre. Os \textit{callbacks} são amplamente utilizados para lidar com operações assíncronas, como chamadas de API, leitura de arquivos ou requisições HTTP. Um exemplo comum é o uso do método \texttt{setTimeout}, que executa uma função de \textit{callback} após um determinado tempo de espera.

Como exemplificado na \autoref{lst:callback-js}, há a função \texttt{fetchData} que simula uma operação assíncrona usando \texttt{setTimeout}. Ela recebe uma função de \textit{callback} como argumento e, quando a operação é concluída, chama a função de \textit{callback} passando o resultado como parâmetro. A função \texttt{processarDados} é passada como \textit{callback} e é responsável por processar os dados obtidos ou lidar com erros, caso ocorram.

\begin{lstlisting}[language=TypeScript, caption={Exemplo de uso de \textit{callback} em \textit{JavaScript}}, label=lst:callback-js]
function fetchData(callback) {
    // Simulando uma operacao assincrona
    setTimeout(() => {
        const data = 'Dados obtidos com sucesso!';
        // Chamando a funcao de callback
        callback(null, data); 
    }, 2000);
}
    
function processarDados(error, data) {
    if (error) {
        console.error('Ocorreu um erro:', error);
        return;
    }  
    console.log('Dados processados:', data);
}
// Passando a funcao de callback como argumento
fetchData(processarDados); 
\end{lstlisting}


O \textit{TypeScript}, como um \textit{superset} do \textit{JavaScript}, também oferece suporte total aos \textit{callbacks}. No entanto, com a adição de recursos avançados de tipagem estática, o \textit{TypeScript} permite uma maior segurança e clareza ao trabalhar com \textit{callbacks}. Na \autoref{lst:callback-ts}, adiciona-se as anotações de tipo ao \textit{callback} e aos parâmetros das funções \texttt{fetchData} e \texttt{processarDados}. O \textit{TypeScript} infere e valida os tipos corretos, proporcionando uma verificação de tipo adicional durante o desenvolvimento. Além disso, é possível utilizar tipos mais complexos para os parâmetros e retorno das funções de \textit{callback}, proporcionando uma maior expressividade e segurança ao código.


\begin{lstlisting}[language=TypeScript, caption={Exemplo de uso de \textit{callback} em \textit{TypeScript}}, label=lst:callback-ts]
function fetchData(callback: (error: Error 
| null, data: string) => void) {
    // Simulando uma operacao assincrona
    setTimeout(() => {
        const data = 'Dados obtidos com sucesso!';
        callback(null, data); // Chamando a funcao de callback
    }, 2000);
}

function processarDados(error: Error | null, data: string) {
    if (error) {
        console.error('Ocorreu um erro:', error);
        return;
    }
    console.log('Dados processados:', data);
}
// Passando a funcao de callback como argumento
fetchData(processarDados);
\end{lstlisting}


\textit{Callbacks} são uma forma poderosa de lidar com operações assíncronas e chamadas de retorno em \textit{JavaScript} e \textit{TypeScript}. Eles permitem que você controle o fluxo de execução e lide com o resultado de uma operação de forma assíncrona. No entanto, o uso excessivo de \textit{callbacks} pode levar a um código complexo e de difícil manutenção, conhecido como \textit{callback hell}. Para evitar o \textit{callback hell} e melhorar a legibilidade do código, outras abordagens, como \textit{Promises} e \texttt{async}/\texttt{await}, foram introduzidas nas versões mais recentes do \textit{JavaScript} e do \textit{TypeScript}. Essas abordagens oferecem uma sintaxe mais limpa e uma melhor estrutura para lidar com operações assíncronas.

\section{Arrays}
\label{JS_arays}

Arrays são estruturas de dados que permitem armazenar e manipular coleções de elementos. Eles podem ser usados para armazenar uma lista de valores, como números, strings, objetos e outros tipos de dados. Para inicializar um array em \textit{JavaScript} e \textit{TypeScript}, usa-se colchetes para delimitar a cadeia de dados pertencente à estrutura. a \autoref{lst:dec-arr} evidencia a sintaxe de declaração de um array e exemplifica arrays de diferentes tipos de elementos. Na segunda linha é declarado um \textit{array} de inteiros \texttt{numeros}, na quarta linha é declarado um \textit{array} de strings \texttt{nomes} e na sexta linha é declarado um \textit{array} vazio \texttt{arr\_vazio}.

\Needspace{0.3\textheight}
\begin{lstlisting}[language=TypeScript, caption=Sintaxe de Inicialização de Array em JS e TS.,label=lst:dec-arr]
// array de inteiros
let numeros = [1, 2, 3, 4, 5];
// array de strings
let nomes = ["Pedro", "Maria", "Felipe"];
// array vazio
let arr_vazio = [];
\end{lstlisting}

Como já mencionado, um array se trata de uma cadeia de elementos, portanto, não é possível acessar um valor específico simplesmente chamando o nome da variável que o armazena. Para acessar um elemento de um array utiliza-se o índice, que é a posição de um elemento dentro de um array. Os índices de um array em \textit{JavaScript} e \textit{TypeScript} começam em zero, ou seja, zero é a primeira posição do array. Na \autoref{lst:arr-index} é possível observar essa indexação de um array. No exemplo é definido um array chamado \texttt{numeros}, o índice 0 desse array é seu primeiro valor, que é o elemento 1, o índice 2 se refere ao seu terceiro elemento, que é 3, e o índice 4 do array indica seu último índice no exemplo, o quinto elemento, que é o número 5.

\begin{lstlisting}[language=TypeScript, caption=Exemplo de Indexação de Array em JS e TS., label=lst:arr-index]
let numeros = [1, 2, 3, 4, 5];
console.log(numeros[0]); // imprime 1
console.log(numeros[2]); // imprime 3
console.log(numeros[4]); // imprime 5
\end{lstlisting}


\section{Coleções Chaveadas}
\label{JS_colecoesChaveadas}

As coleções chaveadas são estruturas de dados que permitem armazenar e manipular conjuntos de elementos em \textit{JavaScript} e \textit{TypeScript} de forma eficiente. Duas das principais coleções chaveadas disponíveis nessas linguagens são os Mapas (\textit{Map}) e os Conjuntos (\textit{Set}).

\textbf{Mapas (\textit{Map}):} Os Mapas são estruturas que permitem armazenar pares de chave-valor, onde cada chave é exclusiva e associada a um valor correspondente. Isso significa que os Mapas permitem mapear valores a partir de uma chave específica. Em \textit{JavaScript} e \textit{TypeScript}, a classe \texttt{Map} é usada para criar e manipular Mapas.

Para criar um Mapa vazio, pode-se usar o construtor da classe Map. Em seguida, é possível adicionar elementos ao Mapa usando o método \texttt{set(chave, valor)}, onde a chave representa a chave do elemento e o valor é o valor associado a essa chave. Na \autoref{lst:map-example}, um Mapa é criado e dois elementos são adicionados a ele. A chave ``chave1'' está associada ao valor ``valor1'', e a chave ``chave2'' está associada ao valor ``valor2''.

\Needspace{0.3\textheight}
\begin{lstlisting}[language=TypeScript, caption=Exemplo de utilização de Mapa em JS e TS., label=lst:map-example]
const mapa = new Map();
mapa.set('chave1', 'valor1');
mapa.set('chave2', 'valor2');
\end{lstlisting}

Para acessar um valor específico no Mapa, pode-se usar o método \texttt{get(chave)}, passando a chave desejada como parâmetro. Para exemplificar, na \autoref{lst:map-access}, é feito o acesso ao valor associado à chave ``chave1'', resultando na exibição do valor ``valor1'' no console, ainda como base o exemplo da \autoref{lst:map-example}. Além disso, os Mapas também possuem métodos para verificar a presença de uma chave (\texttt{has(chave)}), remover um elemento (\texttt{delete(chave)}) e obter a quantidade de elementos no Mapa (\texttt{size}). Essas funcionalidades tornam os Mapas extremamente úteis para armazenar e recuperar informações com base em chaves únicas.

\begin{lstlisting}[language=TypeScript, caption=Acessando um valor específico em um Mapa., label=lst:map-access]
console.log(mapa.get('chave1')); // Saida: 'valor1'
\end{lstlisting}



\textbf{Conjuntos (\textit{Set}):} Os Conjuntos são estruturas que permitem armazenar elementos únicos, sem a necessidade de chaves associadas. Em outras palavras, um Conjunto contém apenas valores distintos, eliminando automaticamente duplicatas. Em \textit{JavaScript} e \textit{TypeScript}, a classe \texttt{Set} é usada para criar e manipular Conjuntos.

Na \autoref{lst:set-example}, um Conjunto é criado a partir da classe \texttt{Set} com o comando ``\texttt{new Set()''} e o Conjunto é atribuído à variável \texttt{conjunto}. Dois elementos são adicionados a ele, ``elemento1'' e ``elemento2'', com o comando \texttt{add(elemento)}. O Conjunto garante que apenas valores únicos sejam armazenados, portanto, se um elemento duplicado for adicionado, ele será ignorado.

\begin{lstlisting}[language=TypeScript, caption=Exemplo de utilização de Conjunto em JS e TS., label=lst:set-example]
const conjunto = new Set();
conjunto.add('elemento1');
conjunto.add('elemento2');
\end{lstlisting}


Na \autoref{lst:set-duplicates}, o elemento ``elemento1'' é adicionado duas vezes ao Conjunto. No entanto, como os Conjuntos não permitem elementos duplicados, apenas uma ocorrência desse elemento será mantida no Conjunto, resultando em um tamanho de 2. Ademais, os Conjuntos também possuem métodos úteis para verificar a existência de um elemento (\texttt{has(elemento)}), remover um elemento (\texttt{delete(elemento)}) e obter a quantidade de elementos no Conjunto (\texttt{size}).
\Needspace{0.3\textheight}

\begin{lstlisting}[language=TypeScript, caption=Removendo elementos duplicados em um Conjunto., label=lst:set-duplicates]
const conjunto = new Set();
conjunto.add('elemento1');
conjunto.add('elemento2');
conjunto.add('elemento1'); 
// Elemento duplicado eh removido automaticamente
console.log(conjunto.size); // Saida: 2
\end{lstlisting}

Em resumo, os Mapas e os Conjuntos são poderosas coleções chaveadas que permitem armazenar e manipular dados de forma eficiente em \textit{JavaScript} e \textit{TypeScript}. Os Mapas são ideais quando é necessário mapear valores a partir de chaves únicas, enquanto os Conjuntos são adequados para armazenar elementos únicos sem a necessidade de associação chave-valor.
%% https://developer.mozilla.org/pt-BR/docs/Web/JavaScript/Guide/Keyed_collections



\section{Classes e Objetos}
\label{JS_classesObjetos}

As classes e objetos são conceitos fundamentais na programação orientada a objetos (POO). Em \textit{JavaScript} e \textit{TypeScript}, é possível criar e manipular classes para representar entidades e estruturar o código de forma mais organizada e reutilizável. Uma classe é uma estrutura que define um conjunto de atributos e métodos relacionados. Os atributos representam as características dos objetos da classe, enquanto os métodos são as ações que esses objetos podem executar. Em \textit{JavaScript} e \textit{TypeScript}, os atributos e métodos são definidos dentro da declaração da classe.

Na \autoref{lst:classe} é possível observar a criação de uma classe \texttt{Carro} em \textit{JavaScript} e \textit{TypeScript}. A classe \texttt{Carro} possui dois atributos: \texttt{cor}, do tipo \texttt{string}, e \texttt{rodas}, do tipo \texttt{number}. Além disso, ela possui um método \texttt{informacoes()} que imprime uma mensagem contendo a cor e a quantidade de rodas do carro. Isso quer dizer que a classe \texttt{Carro} encapsula as propriedades e comportamentos que serão passados às instâncias (objetos) dessa classe.

\Needspace{0.3\textheight}
\begin{lstlisting}[language=TypeScript, caption=Exemplo básico de classe em JS e TS., label=lst:classe]
class Carro {
    cor = "vermelho";
    rodas = 4;
    informacoes() {
        console.log(`Sou um carro ${this.cor}
        e tenho ${this.rodas} rodas.`);
    }
}
\end{lstlisting}

É importante destacar que o \textit{TypeScript} oferece suporte a tipos de dados estáticos, permitindo que você especifique os tipos dos atributos e parâmetros de métodos. Por exemplo, no \textit{TypeScript}, dessa forma, é garantido que a propriedade \texttt{cor} seja sempre uma \texttt{string} e a propriedade \texttt{rodas} seja sempre do tipo \texttt{number}. Essa verificação estática de tipos proporcionada pelo \textit{TypeScript} ajuda a evitar erros e a facilitar a compreensão e manutenção do código.

A instanciação de objetos é o processo de criar uma instância ou um objeto específico de uma classe. Em \textit{JavaScript} e \textit{TypeScript}, isso é feito usando o operador \texttt{new}, seguido do nome da classe e de quaisquer argumentos necessários pelo construtor. Como ilustra a \autoref{lst:instanciacao}, é criada uma instância da classe \texttt{Carro} usando o operador \texttt{new} e essa instância é atribuída à variável \texttt{meuCarro}. Em seguida, o método \texttt{informacoes()} é chamado na instância \texttt{meuCarro} para exibir uma mensagem contendo a cor e a quantidade de rodas do carro.

\begin{lstlisting}[language=TypeScript, caption=Instanciação de objetos em JS e TS., label=lst:instanciacao]
const meuCarro = new Carro();
meuCarro.informacoes();
// Saida: "Sou um carro vermelho e 
//tenho 4 rodas."
\end{lstlisting}


A instanciação de objetos permite criar múltiplas instâncias da mesma classe, cada uma com seus próprios valores de atributos. Isso é útil quando se deseja trabalhar com diferentes objetos que compartilham a mesma estrutura e comportamento definidos pela classe. Por isso, ao criar uma classe, a fim de torná-la versátil pode-se definir uma função construtora para ela, que pode depender de parâmetros, ou não. Em \textit{JavaScript} e \textit{TypeScript}, isso pode ser feito usando o método especial \texttt{constructor()}. Esse método é executado automaticamente quando um objeto da classe é criado.



Na \autoref{lst:constructor-js}, o método \texttt{constructor()} recebe dois parâmetros: cor e rodas. Dentro do inicializador, esses valores são atribuídos aos atributos correspondentes do objeto usando a palavra-chave \texttt{this}. Agora, ao criar um objeto Carro, é necessário fornecer os valores para \texttt{cor} e \texttt{idade} durante a instanciação. A partir disso, cada instância de \texttt{Carro} terá sua própria cor e quantidade de rodas, permitindo a criação e manipulação de vários carros independentes uns dos outros.


\begin{lstlisting}[language=TypeScript, caption=Adicionando inicializador de objetos em JS., label=lst:constructor-js]
class Carro {
    constructor(cor, rodas) {
        this.cor = cor;
        this.rodas = rodas;
    }
    informacoes() {
        console.log(`Sou um carro ${this.cor}
        e tenho ${this.rodas} rodas.`);
    }
}
\end{lstlisting}



O \textit{TypeScript}, sendo uma linguagem de programação baseada em \textit{JavaScript}, adiciona a verificação estática de tipos durante a compilação, garantindo que os valores fornecidos durante a instanciação estejam corretamente tipados de acordo com a definição da classe. a \autoref{lst:constructor-ts} mostra o mesmo exemplo da \autoref{lst:constructor-js}, porém, adaptado para \textit{TypeScript}, explicitando os tipos \texttt{string} e \texttt{number} para os atributos \texttt{cor} e \texttt{rodas}, respectivamente, além de definir também o tipo \texttt{void} para o retorno do método \texttt{informacoes}.


\begin{lstlisting}[language=TypeScript, caption=Adicionando inicializador de objetos em TS., label=lst:constructor-ts]
class Carro {
    cor: string;
    rodas: number;
    constructor(cor: string, rodas: number) {
        this.cor = cor;
        this.rodas = rodas;
    }
    informacoes(): void {
        console.log(`Sou um carro ${this.cor} 
        e tenho ${this.rodas} rodas.`);
    }
}
\end{lstlisting}


Em \textit{JavaScript}, é comum usar uma função construtora para criar objetos de uma classe. Uma função construtora é uma função regular que é invocada usando o operador new e retorna um objeto. Essa abordagem é útil quando se deseja definir membros privados para a classe. Na \autoref{lst:construtora}, a função \texttt{Pessoa} serve como função construtora. Dentro dela, define-se os atributos \texttt{nome} e \texttt{idade} como propriedades do objeto usando \texttt{this}. Também é definido o método \texttt{falar()} como uma função do objeto. Ao criar um novo objeto \texttt{Pessoa} usando o operador \texttt{new}, a função construtora é invocada e retorna um objeto com as propriedades e métodos definidos. Essa maneira também é válida em \textit{TypeScript}, no entanto, o uso da sintaxe de classe com a palavra-chave \texttt{class} é recomendado para aproveitar os recursos adicionais de verificação de tipo e recursos avançados de POO oferecidos pelo \textit{TypeScript}.



\begin{lstlisting}[language=TypeScript, caption=Exemplo de função construtora em JS., label=lst:construtora]
function Pessoa(nome, idade) {
    this.nome = nome;
    this.idade = idade;
    this.falar = function() {
    console.log(`Ola, meu nome eh ${this.nome} 
    e tenho ${this.idade} anos.`);
    };
}
const pessoa = new Pessoa("Joao", 25);
pessoa.falar(); 
// Saida: Ola, meu nome eh Joao e tenho 25 anos.
\end{lstlisting}

Além dessas formas mencionadas, há também o método \texttt{Object.create()}, que é outra forma de criar objetos a partir de uma classe em \textit{JavaScript} e \textit{TypeScript}. Ele permite criar um novo objeto com um protótipo especificado. O protótipo é um objeto existente que será usado como base para o novo objeto. Por exemplo, na \autoref{lst:objectcreate} é definido o objeto \texttt{pessoaProto} que contém o método \texttt{falar()}. Ao chamar \texttt{Object.create(pessoaProto)}, cria-se um novo objeto pessoa que tem \texttt{pessoaProto} como protótipo. Em seguida, os atributos nome e idade no objeto pessoa são definidos. Dessa forma, o objeto pessoa herda o método \texttt{falar()} do protótipo.



\begin{lstlisting}[language=TypeScript, caption=Utilizando \texttt{Object.create} para criar objetos em JS., label=lst:objectcreate]
const pessoaProto = {
    falar: function() {
        console.log(`Ola, meu nome eh ${this.nome} 
        e tenho ${this.idade} anos.`);
    }
};

const pessoa = Object.create(pessoaProto);
pessoa.nome = "Maria";
pessoa.idade = 30;
pessoa.falar(); 
// Saida: Ola, meu nome eh Maria e tenho 30 anos.
\end{lstlisting}


Em \textit{JavaScript}, não há um suporte nativo para membros privados em classes. No entanto, é possível simular membros privados usando convenções de nomenclatura ou recursos avançados do \textit{TypeScript}, como modificadores de acesso. Na \autoref{lst:privado}, os atributos \texttt{nome} e \texttt{idade} são prefixados com \texttt{\#}, indicando que eles são membros privados da classe.  Essa convenção de nomenclatura é uma prática comum para indicar que esses atributos devem ser tratados como privados e não devem ser acessados diretamente fora da classe. Em \textit{JavaScript} pode ser usado apenas como convenção, e não impede de fato o acesso a esses elementos, diferentemente do \textit{TypeScript}, que emitirá um erro ao tentar acessá-los.



\begin{lstlisting}[language=TypeScript, caption=Simulação de membros privados usando convenção de nomenclatura em JS., label=lst:privado]
class Pessoa {
    #nome: string;
    #idade: number;
    constructor(nome: string, idade: number) {
        this.#nome = nome;
        this.#idade = idade;
    }
    falar(): void {
        console.log(`Ola, meu nome eh ${this.#nome}
        e tenho ${this.#idade} anos.`);
    }
}
let a = new Pessoa("Joana", 22);
console.log(a.#nome);
// Saida em JS: "Joana"
/* Saida em TS: "Error: Property '#nome' is not 
accessible outside class 'Pessoa' because it has a
private identifier." */
\end{lstlisting}

O \textit{TypeScript} também oferece modificadores de acesso como o \texttt{public}, \texttt{private} e \texttt{protected} para controlar o acesso aos membros da classe. Esses modificadores fornecem um nível adicional de segurança e encapsulamento em relação ao \textit{JavaScript} puro.  A \autoref{tab:modif-acess} relaciona os três modificadores de acesso em \textit{TypeScript} com suas respectivas usabilidades, dentre eles, o \texttt{private}, que deixa o membro acessível somente dentro da própria classe. a \autoref{lst:private} mostra o mesmo exemplo da \autoref{lst:privado}, utilizando o comando \texttt{private} do \textit{TypeScript}.


\begin{table}
\centering
\begin{tabular}{|c|p{0.6\linewidth}|}
\hline
\textbf{Modificador de Acesso} & \textbf{Significado} \\
\hline
\texttt{private} & O membro (propriedade ou método) é acessível apenas dentro da própria classe em que foi declarado. Não é possível acessá-lo fora da classe ou em classes derivadas. \\
\hline
\texttt{public} & O membro é acessível em qualquer lugar, ou seja, em qualquer classe ou parte do código. \\
\hline
\texttt{protected} & O membro é acessível dentro da própria classe em que foi declarado e em classes derivadas dessa classe. Não é possível acessá-lo fora desses contextos. \\
\hline
\end{tabular}
\caption{\label{tab:modif-acess}Modificadores de Acesso em \textit{TypeScript}}
\end{table}


\Needspace{0.3\textheight}
\begin{lstlisting}[language=TypeScript, caption=Propriedades privadas em TS., label=lst:private]
class Pessoa {
    private nome: string;
    private idade: number;
    constructor(nome: string, idade: number) {
        this.nome = nome;
        this.idade = idade;
    }
    falar(): void {
        console.log(`Ola, meu nome eh ${this.nome} 
        e tenho ${this.idade} anos.`);
    }
}
let a = new Pessoa("Joana", 22);
console.log(a.nome);
/* Error: Property 'nome' is private and only 
accessible within class 'Pessoa'. */
\end{lstlisting}
 

Os membros estáticos em uma classe são aqueles que pertencem à classe em si, em vez de pertencerem a instâncias individuais da classe. Em \textit{JavaScript} e \textit{TypeScript}, os membros estáticos são definidos usando a palavra-chave \texttt{static}. Os membros estáticos são úteis quando se deseja compartilhar informações ou comportamentos que se aplicam à classe como um todo, em vez de instâncias individuais. Por exemplo, é possível utilizar um membro estático para armazenar configurações globais ou criar métodos utilitários que não dependem de um objeto específico. 

Na \autoref{lst:estatico-js}, é apresentado um exemplo de como adicionar um membro estático em \textit{JavaScript}. O membro estático \texttt{pais} é adicionado à classe \texttt{Pessoa}, representando o país padrão de uma pessoa. O membro estático é acessado usando o nome da classe, ou seja, \texttt{Pessoa.pais}. Além disso, é definido um método estático \texttt{mudarPais()}, que permite alterar o país padrão. A classe \texttt{Pessoa} também possui uma propriedade \texttt{nome} e uma propriedade \texttt{idade}, que são específicas de cada instância de pessoa. Porém, o membro \texttt{pais} é compartilhado por todas as instâncias, pois é estático. Isso significa que todas as instâncias de \texttt{Pessoa} terão acesso ao mesmo valor de \texttt{pais}. O método \texttt{falar()} imprime uma mensagem com o nome, idade e país da pessoa, utilizando o membro estático \texttt{pais}. No caso do \textit{TypeScript}, a sintaxe correta deve incluir os tipos dos parâmetros e dos membros da classe explicitamente. 

\begin{lstlisting}[language=TypeScript, caption=Adicionando membro estático em \textit{JavaScript}., label=lst:estatico-js]
class Pessoa {
    nome;
    idade;
    static pais = "Brasil";
    constructor(nome, idade) {
        this.nome = nome;
        this.idade = idade;
    }
    falar() {
        console.log(`Ola, meu nome eh ${this.nome} 
        e tenho ${this.idade} anos. Sou do(a) ${Pessoa.pais}.`);
    }
    static mudarPais(novoPais){
        Pessoa.pais = novoPais;
    }
}
\end{lstlisting}


É importante saber que é possível acessar um membro estático sem a necessidade de criar uma instância da classe. Usando ainda o exemplo da \autoref{lst:estatico-js}, para mudar o valor do país padrão, pode-se chamar o método estático \texttt{mudarPais()} diretamente na classe, utilizando \texttt{Pessoa.mudarPais("Novo País")}. Isso afetará todas as instâncias de \texttt{Pessoa}, as já criadas e as que serão criadas.



\subsection{Herança de Classes}

A herança é um conceito importante na POO que permite criar novas classes com base em classes existentes. A classe que serve como base para a criação de novas classes é chamada de classe pai ou superclasse, e as classes criadas a partir dela são chamadas de classes filhas ou subclasses. Na \autoref{lst:heranca-js}, a classe \texttt{Estudante} herda da classe \texttt{Pessoa}, definida na \autoref{lst:estatico-js}, usando a palavra-chave \texttt{extends}. A classe filha possui um novo atributo matricula e um novo método \texttt{estudar()}. No construtor da classe filha, chama-se o construtor da classe pai usando \texttt{super(nome, idade)} para garantir que os atributos da classe pai sejam corretamente inicializados. Para \textit{TypeScript} a lógica é a mesma, contudo, como já enfatizado, deve-se realizar a especificação explícita dos tipos para os métodos e atributos.

\Needspace{0.3\textheight}
\begin{lstlisting}[language=TypeScript, caption=Herança de classes em \textit{JavaScript}, label=lst:heranca-js]
class Estudante extends Pessoa {
    constructor(nome, idade, matricula) {
        super(nome, idade);
        this.matricula = matricula;
    }
    
    estudar() {
        console.log(`${this.nome} 
        esta estudando. Matricula: ${this.matricula}`);
    }
}
\end{lstlisting}


A herança de classes permite reutilizar a estrutura e o comportamento definidos na classe pai, evitando a duplicação de código. Além disso, as classes filhas podem adicionar novos atributos e métodos específicos ou substituir os existentes, personalizando o comportamento conforme necessário. Também é importante ressaltar que o \textit{TypeScript} fornece um suporte mais robusto para herança de classes, incluindo a verificação de tipos durante a compilação e a capacidade de definir modificadores de acesso para membros herdados.


\section{Interfaces}
\label{JS_interfaces}
As interfaces permitem definir contratos ou estruturas para objetos e fornecem uma forma de garantir que os objetos tenham certos métodos ou propriedades. Em \textit{JavaScript}, as interfaces não são uma construção nativa da linguagem, mas podem ser simuladas usando comentários ou convenções de nomenclatura, como mostra a \autoref{lst:interface-js}. No exemplo, é criada uma função vazia chamada \texttt{PessoaInterface} e são adicionadas propriedades e métodos a ela usando o prototype. Essa função pode ser usada como uma interface para objetos que devem ter propriedades \texttt{nome} e \texttt{idade}, bem como o método \texttt{saudacao()}. É importante observar que essa é apenas uma convenção e não há um mecanismo nativo para verificar a implementação de uma interface em \textit{JavaScript}.

\begin{lstlisting}[language=TypeScript, caption=Simulação de Interface em JS., label=lst:interface-js]
// Definindo uma interface para um objeto 'Pessoa'
function PessoaInterface() {}
PessoaInterface.prototype.nome = '';
PessoaInterface.prototype.idade = 0;
PessoaInterface.prototype.saudacao = function() {};
\end{lstlisting}

Em contraste, o \textit{TypeScript} tem suporte nativo para interfaces e fornece verificação estática de tipos durante a compilação. As interfaces no \textit{TypeScript} são usadas para definir a forma ou estrutura de um objeto. Como é possível observar na \autoref{lst:interface-ts}, é contruída uma interface chamada Pessoa que especifica que um objeto do tipo Pessoa deve ter uma propriedade \texttt{nome} do tipo \texttt{string}, uma propriedade \texttt{idade} do tipo \texttt{number} e um método \texttt{saudacao()} que não retorna nenhum valor (\texttt{void}).

\begin{lstlisting}[language=TypeScript, caption=Exemplo de Interface em TS., label=lst:interface-ts]
// Definindo uma interface para um objeto 'Pessoa'
interface Pessoa {
    nome: string;
    idade: number;
    saudacao(): void;
}
\end{lstlisting}



Pode-se utilizar interfaces para garantir que objetos tenham a estrutura desejada. Na \autoref{lst:uso-interface} há um exemplo de como usar a interface \texttt{Pessoa} em \textit{TypeScript}. No exemplo, a função saudar recebe um parâmetro do tipo \texttt{Pessoa} e imprime uma saudação com base nas propriedades desse objeto. Em seguida foi criada uma variável \texttt{pessoa1} que segue a estrutura definida pela interface \texttt{Pessoa} e esse objeto foi passado para a função \texttt{saudar} e chama-se o método \texttt{saudacao()}. Nesse contexto, o \textit{TypeScript} verifica se o objeto \texttt{pessoa1} cumpre os requisitos da interface \texttt{Pessoa}. Ao adicionar ou remover alguma propriedade ou método necessários pela interface, o \textit{TypeScript} informará sobre o erro durante a compilação.



\begin{lstlisting}[language=TypeScript, caption=Utilização de Interfaces em TS., label=lst:uso-interface]
function saudar(pessoa: Pessoa) {
    console.log(`Ola, meu nome eh ${pessoa.nome} 
    e tenho ${pessoa.idade} anos.`);
}
let pessoa1: Pessoa = {
    nome: 'Joao',
    idade: 30,
    saudacao: function() {
        console.log('Ola, tudo bem?');
    }
};
saudar(pessoa1); 
// Saida: "Ola, meu nome eh Joao e tenho 30 anos."
pessoa1.saudacao(); 
// Saida: "Ola, tudo bem?"
\end{lstlisting}


Violar uma interface em \textit{TypeScript} resultará em erros. Na \autoref{lst:interface-violation-type-error}, a interface \texttt{Animal} é definida com as propriedades \texttt{nome} (do tipo \texttt{string}) e \texttt{idade} (do tipo \texttt{number}). No entanto, ao criar o objeto \texttt{gato}, o valor atribuído à propriedade \texttt{idade} é uma \texttt{string} ao invés de um número. Nesse caso, o \textit{TypeScript} detectará o erro durante a compilação e exibirá uma mensagem de erro indicando a incompatibilidade de tipos.



\begin{lstlisting}[language=TypeScript, caption={Exemplo de erro de tipo ao violar uma interface}, label=lst:interface-violation-type-error]
interface Animal {
    nome: string;
    idade: number;
}
let gato: Animal = {
    nome: 'Bella',
    idade: '2 anos'  
    // Tipo incorreto para a propriedade 'idade'
};
\end{lstlisting}


Ademais, ao tentar acessar uma propriedade ou chamar um método que não esteja definido na interface, o \textit{TypeScript} também nos alertará sobre o erro. Na \autoref{lst:interface-violation-property-method-error}, a interface \texttt{Pessoa} define as propriedades \texttt{nome} e \texttt{idade}, bem como o método \texttt{saudacao()}. No entanto, ao criar o objeto \texttt{pessoa}, adiciona-se um método chamado \texttt{cumprimento} que não está definido na interface. O \textit{TypeScript} exibirá um erro durante a compilação, informando que a propriedade \texttt{cumprimento} não é válida.



\begin{lstlisting}[language=TypeScript, caption={Exemplo de erro ao acessar propriedade e método não definidos na interface}, label=lst:interface-violation-property-method-error]
interface Pessoa {
    nome: string;
    idade: number;
    saudacao(): void;
}
let pessoa: Pessoa = {
    nome: 'Joao',
    idade: 30,
    // Metodo nao definido na interface
    cumprimento: function() {  
        console.log('Ola!');
    }
};
pessoa.cumprimento();  
// Erro: 'cumprimento' nao eh uma funcao valida
\end{lstlisting}


Os exemplos mencionados evidenciam a relevância das verificações de tipo fornecidas pelo \textit{TypeScript} para garantir a conformidade das interfaces. Ao realizar verificações de tipo durante a compilação, o \textit{TypeScript} auxilia na detecção antecipada de erros, contribuindo para prevenir problemas e aprimorar a segurança e a robustez do código. Essa abordagem estática de verificação de tipos permite identificar inconsistências e erros de atribuição de valores em tempo de compilação, o que evita falhas e comportamentos inesperados durante a execução do programa.


\section{Aplicação e Contraste entre \textit{JavaScript} e \textit{TypeScript}}
\label{JS_applicacaoExemplo}
A fim de utilizar as linguagens duas linguagens, \textit{JavaScript} e \textit{TypeScript}, em uma aplicação para destacar os principais contrastes entre ambas, um jogo inspirado no jogo Flappy Bird é usado como exemplo nesta seção.

Na \autoref{lst:jogo-html} tem-se o código \textit{HTML} básico do jogo Flappy Bird. O jogo é definido dentro das tags \texttt{<html>} e \texttt{</html>}. No cabeçalho (\texttt{<head>}), tem-se a inclusão de um arquivo \textit{JavaScript} chamado ``script.js'' usando a tag \texttt{<script>}. Esse arquivo contém o código \textit{JavaScript} responsável pela lógica do jogo. Também temos a definição do título da página como ``Flappy Bird'' e a inclusão de um arquivo de estilo chamado ``style.css'' usando a tag \texttt{<link>}. No corpo (\texttt{<body>}), há dois elementos \texttt{<div>}: um com o id ``bird'' que representa o pássaro do jogo e outro com o id ``score'' que exibe a pontuação do jogador.

\Needspace{0.3\textheight}
\begin{lstlisting}[language=html, caption={HTML do jogo}, label=lst:jogo-html]
<!DOCTYPE html>
<html>
<head>
    <script src="script.js" defer></script>
    <title>Flappy Bird</title>
    <link rel="stylesheet" href="style.css">
</head>
<body>
    <div id="bird"></div>
    <div id="score">Score: 0</div>
</body>
</html>
\end{lstlisting}

Na \autoref{lst:jogo-css} há o código CSS responsável pela estilização do jogo. Inicialmente, são definidas algumas propriedades para o elemento \texttt{body} para remover margens e espaçamentos indesejados. Em seguida, há estilização do elemento \texttt{\#bird}, que representa o pássaro do jogo. Ele é posicionado de forma absoluta (\texttt{position:absolute}) e definido com um tamanho de \texttt{50px} por \texttt{50px} (\texttt{width:50px; height:50px}). A cor de fundo do pássaro é amarela (\texttt{background-color:yellow}). Ainda há a estilização da classe \texttt{.pipe}, que representa os canos do jogo. Os canos também são posicionados de forma absoluta e têm uma largura de \texttt{80px} e altura de \texttt{300px}. A cor de fundo dos canos é verde (\texttt{background-color:green}). Por fim, a estilização do elemento \texttt{\#score}, que exibe a pontuação do jogador. Ele é posicionado de forma absoluta, com uma posição no topo esquerdo da tela (\texttt{top:10px; left:10px}). O tamanho da fonte é definido como \texttt{24px} (\texttt{font-size:24px}).


\begin{lstlisting}[language=html, caption={CSS do jogo}, label=lst:jogo-css]
body{
    overflow: hidden;
    padding: 0;
    margin: 0;
}
#bird {
    position: absolute;
    top: 50%;
    left: 50px;
    width: 50px;
    height: 50px;
    background-color: yellow;
}
.pipe {
    position: absolute;
    width: 80px;
    height: 300px;
    background-color: green;
}
#score {
    position: absolute;
    top: 10px;
    left: 10px;
    font-size: 24px;
}
\end{lstlisting}


No trecho de código da \autoref{lst:jogo-js1} observa-se a declaração de várias variáveis globais utilizadas no jogo. Essas variáveis são utilizadas ao longo do código para controlar o funcionamento do jogo. As variáveis globais são:
\begin{itemize}
    \item \texttt{bird}: uma variável que recebe o elemento \textit{HTML} com o id ``bird'' (o pássaro do jogo).

    \item \texttt{score}: uma variável que recebe o elemento \textit{HTML} com o id ``score'' (o elemento que exibe a pontuação do jogador).

    \item \texttt{pipes}: uma matriz vazia para armazenar os canos do jogo.
    
    \item \texttt{gameSpeed}: uma variável que define a velocidade do jogo.
    
    \item \texttt{maxGravity}: uma constante que define a gravidade máxima.
    
    \item \texttt{gravity}: uma variável que define o valor atual da gravidade.
    
    \item \texttt{jumpHeight}: uma constante que define a altura do salto do pássaro.
    
    \item \texttt{scoreCount}: uma variável que representa a contagem da pontuação do jogador.
    
\end{itemize}



% Parte 1
\begin{lstlisting}[language=TypeScript, caption={Projeto exemplo em JS (Parte 1): Declaração de variáveis globais}, label=lst:jogo-js1]
// Selecionando os elementos HTML necessarios
// elemento com id 'bird'
const bird = document.getElementById('bird'); 
// elemento com id 'score'
const score = document.getElementById('score'); 
const pipes = []; // Array para armazenar os canos
let gameSpeed = 5; // Velocidade do jogo
const maxGravity = 10; // Gravidade maxima
let gravity = 10; // Valor da gravidade atual
const jumpHeight = -30; // Altura do salto
let scoreCount = 0; // Contador de pontuacao
\end{lstlisting}


A definição da função pipeMovement() é feita na \autoref{lst:jogo-js2}, responsável por mover os canos do jogo e verificar as colisões com o pássaro. Dentro da função, há um \textit{loop} que itera sobre todos os canos presentes na matriz pipes. Essa função é chamada repetidamente durante a execução do jogo para movimentar os canos e verificar as colisões. Para cada cano, a função realiza as seguintes ações:
\begin{itemize}
    \item Linha 5: Move o cano para a esquerda de acordo com a velocidade do jogo.

    \item Linhas 7 a 10: Verifica se houve colisão entre o pássaro e o cano usando a função \texttt{isColliding(bird, pipe)}. Caso haja colisão, a função \texttt{gameOver()} é chamada para encerrar o jogo.

    \item Linhas 12 a 16: Verifica se o cano já passou pelo pássaro. Isso é feito comparando a posição do cano com a posição do pássaro. Se o cano já passou pelo pássaro e a propriedade \texttt{passed} do cano é falsa, significa que o jogador ganhou um ponto. Nesse caso, a pontuação é incrementada e exibida no elemento score.
\end{itemize}

\newpage
% Parte 2
\begin{lstlisting}[language=TypeScript, caption={Projeto exemplo em JS (Parte 2): Função de movimentação dos canos}, label=lst:jogo-js2]
function pipeMovement(){
  // Loop para mover e verificar colisoes dos canos
  for (const pipe of pipes) {
    // Movimenta os canos na velocidade do jogo
    pipe.style.left = `${pipe.offsetLeft - gameSpeed}px`;
    // Verifica se o passaro colidiu com o cano
    if (isColliding(bird, pipe)) {
      gameOver();
      callback(()=>{return;});
    }
    // Verifica se o cano ja passou pelo passaro
    if (pipe.offsetLeft + pipe.offsetWidth <
    bird.offsetLeft && !pipe.passed) {
      pipe.passed = true;
      scoreCount++; // aumenta a pontuacao
      score.innerHTML = `Score: ${scoreCount / 2}`;
    }
  }
}
\end{lstlisting}

Já na \autoref{lst:jogo-js3} a função \texttt{pipeSpawn()} é definida, responsável por gerar novos canos e remover os canos que já passaram da tela. Essa função é chamada repetidamente durante a execução do jogo para garantir que sempre haja um número mínimo de canos na tela. Dentro da função, são realizadas as seguintes ações:

\begin{itemize}
    \item Linhas 3 a 5: Verifica se o número de canos presentes na matriz \texttt{pipes} é menor que 4. Se for, chama a função \texttt{generatePipes()} para gerar novos canos.

    \item Linhas 7 a 10: Verifica se o primeiro cano da matriz \texttt{pipes} já passou da tela (se a posição do cano mais a largura do cano for menor que zero). Se sim, remove o primeiro cano da matriz usando o método \texttt{shift()}. Esse método remove o primeiro elemento do \textit{array}.
\end{itemize}


% Parte 3
\newpage
\begin{lstlisting}[language=TypeScript, caption={Projeto exemplo em JS (Parte 3): Função de surgimento de canos}, label=lst:jogo-js3]
function pipeSpawn(){
  // Gera 2 novos canos por vez
  if (pipes.length < 2) {
    generatePipes();
  }
  // Remove os canos que ja passaram da tela
  if (pipes.length > 0 && pipes[0].offsetLeft + 
  pipes[0].offsetWidth < 0) {
    pipes.shift();
  }
}
\end{lstlisting}

Uma outra função, \texttt{applyGravity()}, é declara na \autoref{lst:jogo-js4}, responsável por aplicar a gravidade ao pássaro do jogo. Essa função é chamada repetidamente durante a execução do jogo para simular a queda do pássaro. Dentro da função, são realizadas as seguintes ações:
\begin{itemize}
    \item Linhas 3 a 5: Verifica se a gravidade ainda não atingiu o valor máximo. Se não atingiu, incrementa o valor da gravidade em 2 unidades.

    \item Linha 7: Move o pássaro para baixo com base na gravidade.
\end{itemize}



% Parte 4
\begin{lstlisting}[language=TypeScript, caption={Projeto exemplo em JS (Parte 4): Função de aplicação da gravidade}, label=lst:jogo-js4]
function applyGravity(){
  // Aumenta a gravidade ate o maximo
  if (gravity < maxGravity) {
    gravity += 2;
  }
  // Move o passaro para baixo com base na gravidade
  bird.style.top = `${bird.offsetTop + gravity}px`;
}
\end{lstlisting}


Na \autoref{lst:jogo-js5} é possível observar a declaração da função \texttt{verifyEdges()}, responsável por verificar se o pássaro atingiu os limites da tela, sendo chamada repetidamente durante a execução do jogo. A função verifica se a posição vertical do pássaro somada à sua altura é maior que a altura da janela, ou, se a posição vertical do pássaro é menor que zero. Se alguma dessas condições for verdadeira, significa que o pássaro atingiu os limites da tela. E assim, em caso de colisão com os limites da tela, a função \texttt{gameOver()} é chamada para encerrar o jogo.

\Needspace{0.3\textheight}
% Parte 5
\begin{lstlisting}[language=TypeScript, caption={Projeto exemplo em JS (Parte 5): Função de verificação de bordas da tela}, label=lst:jogo-js5]
function verifyEdges(){
  if (bird.offsetTop + bird.offsetHeight > window.innerHeight ||
  bird.offsetTop < 0) {
    gameOver();
    callback(()=>{return;});
  }
}
\end{lstlisting}


Na \autoref{lst:jogo-js6} tem-se a definição da função \texttt{gameLoop()}, que é a função principal responsável por coordenar o funcionamento do jogo. Dentro da função, são chamadas as seguintes funções em sequência:
\begin{itemize}
    \item \texttt{applyGravity()}: para aplicar a gravidade ao pássaro.

    \item \texttt{pipeMovement()}: para mover os canos.

    \item \texttt{pipeSpawn()}: para gerar novos canos.

    \item \texttt{verifyEdges()}: para verificar se o pássaro atingiu os limites da tela.
\end{itemize}

Após a execução dessas funções, a função \texttt{gameLoop()} é chamada novamente usando \texttt{requestAnimationFrame()}. Essa função cria um \textit{loop} de animação que chama a função \texttt{gameLoop()} a cada novo \textit{frame}. Essa função é chamada uma vez para iniciar o jogo e, em seguida, é chamada repetidamente para manter o jogo em execução.


% Parte 6
\begin{lstlisting}[language=TypeScript, caption={Projeto exemplo em JS (Parte 6): Função de \textit{loop} principal do jogo}, label=lst:jogo-js6]
// Funcao recursiva principal do jogo
function gameLoop() {
  applyGravity();
  pipeMovement();  
  pipeSpawn();
  verifyEdges();
  // Chama a funcao gameLoop para o proximo frame
  requestAnimationFrame(gameLoop);
}
\end{lstlisting}

Ainda, na \autoref{lst:jogo-js7}, há a definição da função \texttt{generatePipes()}, responsável por gerar novos canos para o jogo. Essa função é chamada pela função \texttt{pipeSpawn()} para gerar novos canos durante o jogo. Dentro da função, são realizadas as seguintes ações:
\begin{itemize}
    \item Linhas 4 e 5: Gera duas alturas aleatórias para os canos usando a função \texttt{Math.random()}.

    \item Linhas 8 a 14: Cria quatro elementos de cano usando a função \texttt{createPipeElement()}, passando os parâmetros necessários (classe, posição inicial, posição final e altura).

    \item Linhas 17 a 20: Define a propriedade \texttt{passed} como \texttt{false} para todos os canos gerados.

    \item Linhas 23 a 26: Adiciona os elementos de cano ao corpo do documento.

    \item Linha 27: Adiciona os elementos de cano ao \textit{array} \texttt{pipes}.
\end{itemize}



% Parte 7
\begin{lstlisting}[language=TypeScript, caption={Projeto exemplo em JS (Parte 7): Função de geração de canos}, label=lst:jogo-js7]
// Funcao para gerar novos canos
function generatePipes() {
    // Altura aleatoria dos canos
    const pipeHeight = Math.random() * 
    (window.innerHeight - 400) + 50;

    // Cria os elementos de cano superior e inferior
    const pipeTop = createPipeElement('pipe', 0, 
    window.innerWidth, pipeHeight);
    const pipeBottom = createPipeElement('pipe',
    pipeHeight + 300, window.innerWidth,
    window.innerHeight - pipeHeight - 300);

    // Define a propriedade 'passed'
    pipeTop.passed = false;
    pipeBottom.passed = false;

    // Adiciona os canos ao HTML e ao array de canos
    document.body.appendChild(pipeTop);
    document.body.appendChild(pipeBottom);
    pipes.push(pipeTop, pipeBottom);
}
\end{lstlisting}


A função \texttt{createPipeElement()} é evidenciada na \autoref{lst:jogo-js8}, que é uma função responsável por criar um elemento de cano para o jogo. Essa função é chamada pela função \texttt{generatePipes()} para criar os elementos de cano do jogo. Dentro da função, são realizadas as seguintes ações:
\begin{itemize}
    \item Linha 3: Cria um novo elemento de \texttt{div}.

    \item Linha 4: Define a classe do elemento como a classe recebida como parâmetro.

    \item Linha 5 a 7: Define a posição inicial e a posição final do elemento usando os parâmetros \texttt{top}, \texttt{left} e \texttt{height}.

    \item Linha 8: Retorna o elemento criado.
\end{itemize}

\newpage

% Parte 8
\begin{lstlisting}[language=TypeScript, caption={Projeto exemplo em JS (Parte 8): Função de criação de canos}, label=lst:jogo-js8]
// Funcao para criar um elemento de cano
function createPipeElement(className, top, left, height) {
  const pipeElement = document.createElement('div');
  pipeElement.className = className;
  pipeElement.style.top = `${top}px`;
  pipeElement.style.left = `${left}px`;
  pipeElement.style.height = `${height}px`;
  return pipeElement;
}
\end{lstlisting}


Na \autoref{lst:jogo-js9} há a definição da função \texttt{isColliding()}, responsável por verificar se houve colisão entre dois elementos do jogo. Essa função é chamada pela função \texttt{pipeMovement()} para verificar se houve colisão entre o pássaro e os canos. Dentro da função, são realizadas as seguintes ações:
\begin{itemize}
    \item Verifica se a posição do elemento \texttt{element1} (pássaro) em relação ao lado esquerdo é menor do que a posição do elemento \texttt{element1} (cano) em relação ao lado direito.

    \item Verifica se a posição do elemento \texttt{element1} em relação ao lado direito é maior do que a posição do elemento \texttt{element2} em relação ao lado esquerdo.

    \item Verifica se a posição do elemento \texttt{element1} em relação ao topo é menor do que a posição do elemento \texttt{element1} em relação ao topo.
    
    \item Verifica se a posição do elemento \texttt{element1} em relação ao topo é maior do que a posição do elemento \texttt{element1} em relação ao topo.

    \item Se todas essas condições forem verdadeiras, significa que houve colisão entre os elementos.
\end{itemize}


% Parte 9
\begin{lstlisting}[language=TypeScript, caption={Projeto exemplo em JS (Parte 9): Função de verificação de colisão}, label=lst:jogo-js9]
// Funcao que verifica se dois elementos colidiram
function isColliding(element1, element2) {
  const rect1 = element1.getBoundingClientRect();
  const rect2 = element2.getBoundingClientRect();
  return (
    rect1.top < rect2.bottom &&
    rect1.bottom > rect2.top &&
    rect1.left < rect2.right &&
    rect1.right > rect2.left
  );
}
\end{lstlisting}


A definição da função \texttt{gameOver()} é realizada na \autoref{lst:jogo-js10}, que é responsável por encerrar o jogo quando ocorre uma colisão entre o pássaro e um cano. Essa função é chamada pela função \texttt{pipeMovement()} quando ocorre uma colisão entre o pássaro e um cano, encerrando o jogo. Dentro da função, são realizadas as seguintes ações:
\begin{itemize}
    \item Define a velocidade do jogo como zero para parar o movimento dos canos.

    \item Exibe uma mensagem de ``Game Over'' na página usando o método \texttt{innerHTML} do elemento \texttt{score}.

    \item Remove os eventos de clique (\textit{mousedown}) e de pressionar tecla (\textit{keydown}) do documento para impedir a interação do jogador após o fim do jogo.
\end{itemize}



% Parte 10
\begin{lstlisting}[language=TypeScript, caption={Projeto exemplo em JS (Parte 10): Função de fim de jogo}, label=lst:jogo-js10]
// Funcao para encerrar o jogo
function gameOver() {
  alert(`Fim do jogo! Seu Score: ${scoreCount / 2}`);
  location.reload();
}
\end{lstlisting}

A última função declarada no código é a função \texttt{jump()}, que é apresentada na \autoref{lst:jogo-js11}. Essa função realiza uma verificação condicional que checa se o código da tecla pressionada é igual a 32 (que representa a barra de espaço) e se a variável \texttt{gravity} é maior ou igual a -2. Essa condição garante que o pássaro só pode saltar se a barra de espaço for pressionada e se a gravidade permitir o salto. Se a condição for verdadeira, o valor de \texttt{gravity} é incrementado pela variável \texttt{jumpHeight}. Isso significa que o pássaro ganhará uma força extra para subir no jogo, pois a gravidade será temporariamente reduzida ou anulada.

% Parte 11
\begin{lstlisting}[language=TypeScript, caption={Projeto exemplo em JS (Parte 11): Função de pulo}, label=lst:jogo-js11]
// Funcao para lidar com o salto do passaro
function jump(event) {
  // Verifica se a barra de espaco esta pressionada
  // e se a gravidade permite o salto
  if (event.keyCode === 32 && gravity >= -2) {
    gravity += jumpHeight;
  }
}
\end{lstlisting}

Por fim, como mostra a \autoref{lst:jogo-js12}, para configurar o controle do jogo, em escopo global, um gatilho de evento \textit{keydown} (tecla pressionada) é adicionado para chamar a função \texttt{jump}, fazendo com que, ao pressionar a tecla espaço (verificação realizada dentro da função \texttt{jump}, o passáro pule. E logo em seguida, o jogo é inicializado pela função recursiva \texttt{gameLoop()}.

% Parte 12
\begin{lstlisting}[language=TypeScript, caption={Projeto exemplo em JS (Parte 12): Inicialização do jogo e evento de pulo ao pressionar enter}, label=lst:jogo-js12]
// salto ao pressionar enter
document.addEventListener('keydown', jump);
// Inicia o loop do jogo
gameLoop();
\end{lstlisting}

No código \textit{TypeScript} (TS), há algumas diferenças em relação ao código \textit{JavaScript} (JS). a \autoref{lst:jogo-ts} apresenta as mudanças em relação à alguns trecho de código anteriormente apresentados. Uma das principais diferenças é o uso de tipos estáticos e a declaração de interfaces para elementos do DOM. As principais diferenças entre os códigos \textit{JavaScript} e \textit{TypeScript} são: 
\begin{itemize}
    \item No início do código, a interface \texttt{PipeElement} é definida para estender a interface \texttt{HTMLDivElement} e adicionar a propriedade booleana \texttt{passed}, como está evidenciado na segunda a quarta linha do código. 

    \item Definição do tipo da variável \texttt{pipes} como um \textit{array} de elementos que estejam de acordo com a interface \texttt{PipeElement} criada, como pode ser observado na quinta linha.
    
    \item Dentro da função \texttt{generatePipe}, as variávels dos canos gerados (\texttt{pipeTop} e \texttt{pipeBottom}) recebem elementos \textit{HTML} convertidos para o molde definido pela interface \texttt{PipeElement}, adicionando o atributo \texttt{passed}. A sexta à décima primeira linha ilustram essa mudança.

     \item Os parâmetros da função \texttt{createPipelineElement} têm tipos definidos, sendo \texttt{ClassName} do tipo \texttt{string} e \texttt{top}, \texttt{left} e \texttt{height} do \texttt{number}, como pode ser visto na décima segunda linha.
    
    \item Os parâmetros da função \texttt{isColliding} têm tipos definidos como \texttt{HTMLElement}, como pode ser visto na décima terceira linha.
    
    \item Da mesma forma, a função \texttt{jump} tem seu parâmetro com tipo definido \texttt{KeyboardEvent}, como está definido na décima quarta linha.
\end{itemize}

\Needspace{0.3\textheight}
\begin{lstlisting}[language=TypeScript, caption={Mudanças no código do projeto exemplo ao desenvolver em \textit{TypeScript}}, label=lst:jogo-ts]
// Interface para o elemento do cano
interface PipeElement extends HTMLDivElement {
    passed: boolean;
}
const pipes: PipeElement[] = [];
function generatePipes() {
    ...
    const pipeTop = document.createElement('div') as
    PipeElement;
    const pipeBottom = document.createElement('div') as
    PipeElement;
    ...
}
function createPipeElement(className:string, top: number, 
left: number, height: number) {...}
function isColliding(element1: HTMLElement, element2: 
HTMLElement) {...}
function jump(event: KeyboardEvent) {...}
\end{lstlisting}

Além disso, uma mudança não apresentada na \autoref{lst:jogo-ts}, mas que é importante mencionar é que onde aparecem as variáveis \texttt{bird} e \texttt{score} deve-se usar um símbolo de exclamação (!) posterior à variável, por exemplo: \texttt{bird!.offsetTop}. A exclamação após a variável indica asserção de que a variável não terá valor nulo. Essa mudança é necessária no código, pois o \textit{TypeScript} verifica que há a possibilidade da variável ter valor nulo e dá erro.

Ao desenvolver o projeto\footnote{O código desse projeto está disponível no GitHub em https://github.com/JadsonFaustino/flappy-box-game-js-ts.} exemplo proposto nesta seção é possível evidenciar algumas características do \textit{TypeScript} que se sobressaem sobre o \textit{JavaScript} pela perspectiva de desenvolvimento:
\begin{itemize}
    \item As variáveis são inferidas dinamicamente em tempo de execução no \textit{JavaScript}, permitindo que um valor de qualquer tipo seja atribuído a uma variável. Já no \textit{TypeScript}, as variáveis podem ter seus tipos declarados explicitamente, o que traz mais segurança ao código e facilita a detecção de erros durante a fase de desenvolvimento.

    \item Embora \textit{JavaScript} seja uma linguagem de POO, o \textit{TypeScript} aprimora esse aspecto com recursos adicionais. \textit{TypeScript} permite a definição de classes, interfaces, herança e polimorfismo de forma mais estruturada e explícita. Esses recursos facilitam a organização do código, tornando-o mais legível e mantível. Além disso, o TypeScript também suporta módulos, permitindo a criação de um código modular e reutilizável.

    \item Enquanto o \textit{JavaScript} é interpretado pelo navegador ou pelo ambiente de execução, o \textit{TypeScript} requer um processo de compilação antes da execução. Além disso, o \textit{TypeScript} possui um ecossistema robusto, com ferramentas avançadas de desenvolvimento, como autocompletar, refatoração de código, detecção de erros em tempo real e suporte a IDEs populares. Isso torna o desenvolvimento em \textit{TypeScript} mais eficiente e produtivo.
\end{itemize}

